\section{Helmholtz-SDE}\label{app:sec:helm-sde}

In this section, we provide further explanation of our algorithm, Helmholtz-SDE.

First, in \Cref{app:subsec:helmholtz-compute} we describe how the $q$-weighted Helmholtz decomposition can be obtained for arbitrary smooth, strictly positive marginal densities $q$ by inverting a Stein operator.
For the remainder of \Cref{app:sec:helm-sde}, we then specialize to the setting of Gaussian one-time marginals.
In \Cref{app:subsec:helmholtz-gaussian}, we provide sufficient conditions under which the Helmholtz decomposition exists uniquely.
In \Cref{app:subsec:resid-approx}, we translate this theoretical result into a practical formula using an order-$\ell$ polynomial approximation to the optimal divergence-free correction.
In \Cref{app:subsec:verify-assumptions}, we verify that the resulting variational drift, obtained by adding the polynomial correction to the reference drift, produces an SDE satisfying our assumptions in \Cref{app:subsec:assumptions}.
In \Cref{app:subsec:optimality}, we discuss the optimality of Helmholtz-SDE in the conjugate and non-conjugate settings.
In \Cref{app:subsec:alg-summary} we provide pseudocode of our algorithm.

\subsection{Computing the Helmholtz decomposition}\label{app:subsec:helmholtz-compute}

We first discuss how the Helmholtz decomposition can be computed for an arbitrary smooth, strictly positive density $q$ as the inverse of a Stein operator. The resulting probabilistic expression can be approximated via Monte Carlo. For ease of notation, we consider a state-independent diffusion coefficient, although the same construction holds in the state-dependent setting.

\paragraph{The Helmholtz decomposition and diffusion Stein operator.}
As described in \Cref{subsec:helmholtz-decomp}, the optimal correction to the reference drift $\mbf_q(\mbx, t)$ is obtained by computing the $q$-weighted Helmholtz decomposition of the residual vector field $\mbr(\mbx, t) = \mbf_p(\mbx, t) - \mbf_q(\mbx, t)$ and taking the $q$-divergence-free vector field.
In particular, we want to find a decomposition of $\mbr$ as
\begin{equation}\label{app:eqn:helm-decomp}
    \mbr(\mbx, t) = \mbSigma(t) \nabla_{\mbx} \psi(\mbx, t) + \mbh(\mbx, t), \quad (\mbx, t) \in \bbR^K \times [0, T],
\end{equation}
where $\nabla_{\mbx} \psi(\mbx, t)$ is a gradient field and $\mbh(\mbx, t)$ is a $q$-divergence-free vector field.

Suppose $q(\mbx, t)$ is a smooth, strictly positive density.
To obtain such a decomposition, observe that if \cref{app:eqn:helm-decomp} holds, then multiplying both sides by $q(\mbx, t)$ and computing the divergence yields
\begin{equation}
     \nabla_{\mbx} \cdot [q(\mbx, t) \mbSigma(t) \nabla_{\mbx} \psi(\mbx, t)]
     =
     \nabla_{\mbx} \cdot [q(\mbx, t) \mbr(\mbx, t)],
     \quad
     (\mbx, t) \in \bbR^K \times [0, T],
\end{equation}
since $\mbh$ is $q$-divergence-free. 
Dividing both sides by $q$ and expanding the left-hand side yields
\begin{equation}\label{app:eqn:stein-operator-eqn}
    \mathrm{tr}(\mbSigma(t) \nabla_{\mbx}^2 \psi(\mbx, t))
    +
    \nabla_{\mbx} \log q(\mbx, t) \cdot
    \mbSigma(t) \nabla_{\mbx} \psi(\mbx, t)
    =
    \frac{1}{q(\mbx, t)}
    \nabla_{\mbx} \cdot [q(\mbx, t) \mbr(\mbx, t)] .
\end{equation}

From \cref{app:eqn:stein-operator-eqn}, we make two observations. First, the left-hand side can be written as $[\cL_t^{\mbSigma}\nabla_{\mbx}\psi(\cdot,t)](\mbx)$, where $\cL_t^{\mbSigma}$ is the differential operator
\begin{equation}\label{app:eqn:diffusion-stein}
    [\cL_t^{\mbSigma}\mbf](\mbx) = \mathrm{tr}(\mbSigma(t) D_{\mbx} \mbf (\mbx, t))
    +
    \mbSigma(t) \nabla_{\mbx} \log q(\mbx, t) \cdot
    \mbf(\mbx, t)
\end{equation}
defined for differentiable vector fields $\mbf: \bbR^K \to \bbR^K$, where $D_{\mbx} \mbf$ represents the Jacobian of $\mbf$.
This is the \emph{diffusion Stein operator} corresponding to the density $q(\mbx,t)$ and constant covariance coefficient $\mbSigma(t)$ \citep{barbour1990stein,gorham2019measuring}. 
When $\mbSigma(t)=\bbI_K$, it reduces to the \emph{Langevin Stein} operator
\begin{equation*}
    [\cL_t \mbf](\mbx) = \nabla_{\mbx} \cdot \mbf(\mbx)  + \nabla_{\mbx} \log q(\mbx, t) \cdot \mbf(\mbx).
\end{equation*}
We refer the reader to \citet{gorham2019measuring} for background on diffusion Stein operators.

Second, whenever $\mbr(\mbx, t)$ and
$\cL_t^{\mbSigma}[\mbr(\cdot, t)](\mbx)$ are integrable, the right-hand side of \cref{app:eqn:stein-operator-eqn} has expectation
\begin{equation*}
    \bbE_{\bbQ}\left[
    \frac{1}{q(\mbx, t)}
    \nabla_{\mbx} \cdot [q(\mbx, t) \mbr(\mbx, t)]
    \right]
    =
    \int
    \nabla_{\mbx} \cdot [q(\mbx, t) \mbr(\mbx, t)]
    d \mbx = 0
\end{equation*}
The final equality holds by the divergence theorem \citep[][Theorem 3.1]{liu2026probabilistic}. 

Overall, we have reduced computing the Helmholtz decomposition to solving
\begin{equation}\label{app:eqn:helmholtz-stein-eqn}
\begin{aligned}
    &[\cL_t^{\mbSigma}\nabla_{\mbx} \psi(\cdot, t)](\mbx)
    =
    \eta(\mbx, t),
    \quad
    (\mbx, t) \in \bbR^K \times [0, T],\\
    &\eta(\mbx, t)
    =
    \frac{1}{q(\mbx, t)}
    \nabla_{\mbx} \cdot [q(\mbx, t) \mbr(\mbx, t)]
\end{aligned}
\end{equation}
where $\eta(\mbx, t)$ has mean zero under $q(\mbx,t)$.
Intuitively, in order to solve \cref{app:eqn:helmholtz-stein-eqn}, we need to find an inverse of the diffusion Stein operator $\cL_t^{\mbSigma}$. 
Doing so has been the subject of prior literature, namely \citet{barbour1990stein,pardoux2005poisson,gorham2019measuring}.

\paragraph{Inverting the diffusion Stein operator.}

Fix $t \in [0,T]$. Consider the time-homogeneous diffusion in an auxiliary time variable $s$ with diffusion coefficient $\mbSigma(t)^{1/2}$:
\begin{equation}\label{app:eqn:langevin-diffusion}
    d\mbz^{(t)}(s)
    =
    \frac{1}{2} \mbSigma(t)\nabla_{\mbx}\log q(\mbz^{(t)}(s), t)ds
    +
    \mbSigma(t)^{1/2}d\mbw(s). 
\end{equation}
When $\mbSigma(t) = \bbI_K$, \cref{app:eqn:langevin-diffusion} is recognizable as the overdamped Langevin diffusion.
In particular, \cref{app:eqn:langevin-diffusion} has invariant density $q(\mbx, t)$.
Its \emph{infinitesimal generator} is
\begin{equation}
\begin{aligned}
    [\cA_t u](\mbz)
    &=
    \frac{1}{2q(\mbz, t)}
    \nabla_{\mbx}\cdot
    \left[
        q(\mbz, t)\mbSigma(t)\nabla_{\mbz}u(\mbz)
    \right]\\
    &= \frac{1}{2} \mathrm{tr}(\mbSigma(t)\nabla_{\mbz}^2 u(\mbz)) + \frac{1}{2} \mbSigma(t) \nabla_{\mbz}\log q(\mbz, t) \cdot \nabla_{\mbz} u(\mbz)
\end{aligned}
\end{equation}
for $u: \bbR^K \to \bbR$ twice continuously differentiable and satisfying integrability conditions \cite{karatzas2014brownian}.
This implies, by \cref{app:eqn:diffusion-stein}, $\cA_t u = \cL_t^{\mbSigma}(\frac{1}{2}\nabla_{\mbx}u)$. 
Hence, solving \cref{app:eqn:helmholtz-stein-eqn} is equivalent to solving
\begin{equation}\label{app:eqn:poisson}
    [\cA_t u(\cdot, t)](\mbx) = \eta(\mbx, t), \quad  (\mbx, t) \in \bbR^K \times [0, T]
\end{equation}
and then setting $\nabla_{\mbx}\psi(\mbx,t) = \frac{1}{2} \nabla_{\mbx}u(\mbx, t)$.
\Cref{app:eqn:poisson} is known as the \emph{Poisson equation} \cite{pardoux2005poisson,gilbarg1998elliptic}.

Assume that, for each fixed $t$, $\eta(\mbx, t)$ and the coefficients of the above diffusion satisfy the regularity conditions of \citet[][Theorem 5]{gorham2019measuring}. Then \citet[][Theorem 5]{gorham2019measuring} implies that 
\begin{equation}\label{app:eqn:poisson-soln}
    u(\mbx, t) = -\int_0^\infty \bbE[\eta(\mbz^{(t)}(s), t) | \mbz^{(t)}(0) = \mbx]ds
\end{equation}
solves the Poisson equation. The integrand in \cref{app:eqn:poisson-soln} is computed by starting the diffusion \cref{app:eqn:langevin-diffusion} from $\mbx$ and then running it forward until time $s$.

Therefore,
\begin{equation}\label{app:eqn:helmholtz-soln}
    \nabla_{\mbx}\psi(\mbx,t) =  -\frac{1}{2}\nabla_{\mbx} \int_0^\infty \bbE[\eta(\mbz^{(t)}(s), t) | \mbz^{(t)}(0) = \mbx]ds
\end{equation}
solves \cref{app:eqn:helmholtz-stein-eqn}. Consequently, defining $\mbh(\mbx,t) = \mbr(\mbx,t) - \mbSigma(t)\nabla_{\mbx}\psi(\mbx,t)$
gives $\nabla_{\mbx} \cdot \left[ q(\mbx,t)\mbh(\mbx,t) \right] = 0$.
Thus $\mbh(\cdot,t)$ is the $q$-divergence-free component of the residual vector field.

\subsection{The Helmholtz decomposition for Gaussian marginals}\label{app:subsec:helmholtz-gaussian}

For the remainder of \Cref{app:sec:helm-sde}, we specialize to the setting of Gaussian one-time marginals and a state-independent diffusion coefficient.
In this subsection, we establish the existence,  uniqueness, and regularity of the weighted Helmholtz decomposition in this setting.

Since we are considering only the existence and uniqueness of the weighted Helmholtz decomposition, and not the implied variational SDE, we do not require the assumptions from \Cref{app:subsec:assumptions}.

First, we show that under minimal assumptions, the $q$-divergence-free vector field $\mbh$ that minimizes \cref{eqn:optimize} exists. It is the orthogonal projection onto the subspace of $q$-divergence-free vector fields.
\begin{theorem}[Existence and uniqueness of weak Helmholtz decomposition]\label{app:thm:hilbert-projection}
     Suppose that $\mbr(\mbx, t)$ has polynomial growth and that $\mbSigma(t)$ is positive definite. 
     Let $\cL^2(q(\mbx, t), \bbR^K, \mbSigma(t))^K$ be the Hilbert space of vector fields $\mbf: \bbR^K \to \bbR^K$ satisfying $\bbE_{q(\mbx, t)}\|\mbf(\mbx)\|^2 < \infty$ equipped with the inner product 
     \begin{equation}\label{app:eqn:l2-inner-prod}
         \langle \mbf, \mbg \rangle_{q, t} = \bbE_{q(\mbx, t)} [\mbf(\mbx)^\top \mbSigma(t)^{-1} \mbg(\mbx)], \qquad \| \mbf \|_{q, t} = \langle \mbf, \mbf \rangle_{q, t}^{1/2}
    \end{equation}

     \smallskip
     Then $\mbr(\mbx, t) \in \cL^2(q(\mbx, t), \bbR^K, \mbSigma(t))^K$ admits the decomposition 
     \begin{equation}\label{app:eqn:helmholtz-weak}
         \mbr(\mbx, t) = \mbg(\mbx, t) + \mbh(x, t)
     \end{equation}
     where $\mbh, \mbg \in  \cL^2(q(\mbx, t), \bbR^K, \mbSigma(t))^K$, $\mbh$ is $q$-divergence-free in the weak sense \cref{app:eqn:div-free-weak}, and $\mbg$ is orthogonal to all $q$-divergence-free vector fields.
     Moreover, this decomposition is unique in $\cL^2(q(\mbx, t), \bbR^K, \mbSigma(t))^K$. 
\end{theorem}
\begin{proof}
    Consider the set $\cV_t$ of all vector fields that are divergence-free in the weak sense \cref{app:eqn:div-free-weak}:
    \begin{equation}\label{app:eqn:weak-subspace}
    \cV_t = \left\{ \mbv \in  \cL^2(q(\mbx, t), \bbR^K, \mbSigma(t))^K \ \bigg| \  \int \mbv(\mbx, t)^\top \nabla_{\mbx} \varphi(\mbx) q(d\mbx, t) = 0, \quad \forall \varphi \in \cC \right\},
    \end{equation}
    where $\cC$ is the collection of test functions that are infinitely differentiable and have compact support.

    Since $\mbr(\mbx, t)$ has polynomial growth, it belongs to $\cL^2(q(\mbx, t), \bbR^K, \mbSigma(t))^K$.
    By the Hilbert projection theorem, the result follows once we establish that $\cV_t$ is a closed subspace of $\cL^2(q(\mbx, t), \bbR^K, \mbSigma(t))^K$. 
   
    The fact that $\cV_t$ is a subspace is immediate from the definition. To show that it is closed, fix $\varphi \in \cC$ and define
    \begin{equation*}
    \cL_\varphi(\mbv) =  \int \mbv(\mbx)^\top\nabla_{\mbx}\varphi(\mbx) q(d\mbx,t).
    \end{equation*}
Since $\cL_\varphi(\mbv) = \left\langle \mbv, \mbSigma(t)\nabla_{\mbx}\varphi \right\rangle_{q, t}$, the Cauchy--Schwarz inequality gives
\begin{equation*}
|\cL_\varphi(\mbv)| \leq \|\mbv \|_{q, t} \left\| \mbSigma(t)\nabla{\mbx}\varphi \right\|_{q, t}.
\end{equation*}
Consequently, $\cL_{\varphi}$ is a continuous linear functional.
Therefore, $\cV_t = \bigcap_{\varphi\in \cC}
\mathrm{ker} \cL_{\varphi}$ is closed.
\end{proof}

The preceding theorem tells us that, among the vector fields that are $q$-divergence-free in the weak sense \cref{app:eqn:div-free-weak}, $\mbh(\mbx,t)$ uniquely minimizes \cref{eqn:optimize}.
We call \cref{app:eqn:helmholtz-weak} the \emph{weak Helmholtz decomposition}.
\Cref{app:thm:hilbert-projection} does \emph{not}, however, provide a Helmholtz decomposition as in \cref{eqn:helmholtz-decomp}.
Specifically, it neither establishes that $\mbh$ satisfies the $q$-divergence-free condition pointwise, nor identifies $\mbg$ with a gradient field of the form $\mbSigma(t)\nabla_{\mbx}\psi$, nor verifies the regularity of $\mbg$, $\mbh$, or $\psi$.

The following result supplies this representation and regularity for the Gaussian marginals used in Helmholtz-SDE.
 
\begin{assumption}\label{app:assum:resid}
    Assume that the following hold:
    \begin{enumerate}[label=(\roman*)]
        \item\label{item:ass:cont} $(\mbm(t), \mbS(t))_{0 \leq t \leq T}$ is continuously differentiable in $t$, and $\mbS(t)$ is positive definite.
        \item\label{item:ass:diff} $\mbr$ is twice differentiable in $\mbx$, and its derivatives are continuous in $(\mbx, t)$.
        \item\label{item:ass:poly} $\mbr$, $D_{\mbx} \mbr$, and $D_{\mbx}^{(2)}\mbr$ have polynomial growth in $\mbx$, uniformly in $t$.
        \item\label{item:ass:cont-sigma} $\mbSigma(t)$ is continuous and positive definite.
    \end{enumerate}  
    $D_{\mbx} \mbr \in \bbR^{K \times K}$ and $D_{\mbx}^{(2)} \mbr \in \bbR^{K \times K \times K}$ represent the derivative tensors of the residual vector field $\mbr: \bbR^K \to \bbR^K$. 
\end{assumption}

\begin{theorem}[Regularity of the Helmholtz decomposition]\label{app:thm:helm-gaussian}
    Suppose \Cref{app:assum:resid} holds.
    For every $t\in[0,T]$, the weak Helmholtz decomposition of $\mbr(\mbx, t)$ in \cref{app:eqn:helmholtz-weak} has the pointwise representation
    \begin{equation*}
    \mbr(\mbx,t)=\mbSigma(t)\nabla_{\mbx}\psi(\mbx,t)+\mbh(\mbx,t),
    \end{equation*}
    where $\mbh$ is $q$-divergence-free in the strong sense
    \begin{equation*}
    \nabla_{\mbx}\cdot[q(\mbx,t)\mbh(\mbx,t)]=0, \quad (\mbx, t) \in \bbR^K \times [0, T].
\end{equation*}

    $\psi$ is twice differentiable in $\mbx$, $\mbh$ is once differentiable in $\mbx$, and their derivatives are continuous in $(\mbx, t)$.
    Within the stated regularity class, $\mbh$ is unique pointwise, and $\psi(\cdot,t)$ is unique up to an additive constant, which may depend on $t$.
\end{theorem}
\begin{proof}
Fix $t \in [0, T]$.
Since $\mbS(t)$ is positive definite, $q(\mbx, t)$ is a smooth, strictly positive density on $\bbR^K$. 
Recall from \Cref{app:subsec:helmholtz-compute} that solving the Stein operator equation \cref{app:eqn:helmholtz-stein-eqn} for $\nabla_{\mbx}\psi$ reduces, via the identity $\cA_t u = \cL_t^{\mbSigma}(\tfrac{1}{2}\nabla_{\mbx} u)$, to solving the Poisson equation \cref{app:eqn:poisson} for $u$:
\begin{equation}\label{app:eqn:weighted-poisson-theorem}
    [\cA_t u(\cdot,t)](\mbx) \;=\; \frac{1}{2 q(\mbx, t)} \nabla_{\mbx} \cdot \left[q(\mbx, t)\mbSigma(t)\nabla_{\mbx} u(\mbx,t)\right] \;=\; \eta(\mbx, t),
\end{equation}
where $\cA_t$ is the infinitesimal generator of the Langevin diffusion \cref{app:eqn:langevin-diffusion} and $\eta(\mbx, t)$ is the source term in \cref{app:eqn:helmholtz-stein-eqn}.
The gradient component of the Helmholtz decomposition is then recovered as $\nabla_{\mbx}\psi(\mbx, t) = \tfrac{1}{2}\nabla_{\mbx} u(\mbx, t)$.

Our goal is to show that the candidate $u(\mbx, t)$ in \cref{app:eqn:poisson-soln} satisfies \cref{app:eqn:weighted-poisson-theorem} pointwise, and that $u$ is twice continuously differentiable in $\mbx$ with derivatives continuous in $(\mbx, t)$.
Once \cref{app:eqn:weighted-poisson-theorem} is established, defining
\begin{equation}\label{app:eqn:h-definition-theorem}
    \mbh(\mbx,t) = \mbr(\mbx,t) - \mbSigma(t)\nabla_{\mbx}\psi(\mbx,t)
\end{equation}
gives $\nabla_{\mbx} \cdot \left[q(\mbx, t)\mbh(\mbx,t)\right] = 0$, yielding the desired Helmholtz decomposition.

Our proof proceeds in four main steps:
\begin{enumerate}[label=(\roman*)]
    \item\label{item:whiten} Reexpress the candidate $u$ in whitened and rotated coordinates, where the associated Langevin diffusion becomes a product of one-dimensional Ornstein--Uhlenbeck (OU) processes.
    \item\label{item:reg} Prove that the candidate $u$ is well-defined and satisfies the desired regularity conditions.
    \item\label{item:poisson} Verify that $u$ satisfies the Poisson equation \cref{app:eqn:weighted-poisson-theorem} pointwise.
    \item\label{item:unique} Identify the constructed vector fields with the unique weak Helmholtz decomposition from \Cref{app:thm:hilbert-projection}.
\end{enumerate}

\medskip
\textbf{Step \ref{item:whiten}}:
The candidate $u$ proposed in \Cref{app:subsec:helmholtz-compute} (\cref{app:eqn:poisson-soln}) is expressed as an integral of expectations along the Langevin diffusion \cref{app:eqn:langevin-diffusion} with invariant law $q(\mbx, t)$. We use the same whitening and rotation as in \Cref{app:subsec:resid-approx}. Define
\begin{equation*}
\begin{aligned}
    \mbC(t) &= \mbS(t)^{-1/2}\mbSigma(t)\mbS(t)^{-1/2} = \mbQ(t)\mbLambda(t)\mbQ(t)^\top,\\
    \mbLambda(t) &= \diag(\lambda_1(t),\ldots,\lambda_K(t)), \qquad \mbQ(t)^\top\mbQ(t)=\bbI_K,
\end{aligned}
\end{equation*}
and introduce the coordinates
\begin{equation*}
    \tilde\mbx = \mbQ(t)^\top\mbS(t)^{-1/2}(\mbx-\mbm(t)), \qquad \chi(\tilde\mbx,t) = u(\mbm(t)+\mbS(t)^{1/2}\mbQ(t)\tilde\mbx,t).
\end{equation*}
Transform the residual in the same way:
\begin{equation*}
    \widetilde{\mbrho}(\tilde\mbx,t) = \mbQ(t)^\top\mbS(t)^{-1/2} \mbr(\mbm(t)+\mbS(t)^{1/2}\mbQ(t)\tilde\mbx,t).
\end{equation*}
Then $q(\tilde\mbx,t)=\cN(\mb0,\bbI_K)$. Since $\mbS(t)$ and $\mbSigma(t)$ are positive definite, every $\lambda_i(t)$ is strictly positive.

After an orthogonal transformation of the driving Brownian motion, the Langevin diffusion \cref{app:eqn:langevin-diffusion} becomes
\begin{equation}\label{app:eqn:ou-process-theorem}
    d\tilde\mbz^{(t)}(s) = -\tfrac{1}{2}\mbLambda(t)\tilde\mbz^{(t)}(s)ds + \mbLambda(t)^{1/2}d\mbw(s), \qquad \tilde\mbz^{(t)}(0)=\tilde\mbx.
\end{equation}
Thus its coordinates are independent one-dimensional OU processes with invariant law $\cN(0,1)$. Equivalently, each coordinate is a standard OU process after rescaling the artificial time $s$ by $\lambda_i(t)$. The joint invariant law is $\cN(\mb0,\bbI_K)$, and the infinitesimal generator is
\begin{equation*}
    (\tilde{\cA}_t f)(\tilde\mbx) = \tfrac{1}{2}\mathrm{tr}\left( \mbLambda(t)\nabla_{\tilde\mbx}^2 f(\tilde\mbx) \right) - \tfrac{1}{2}\tilde\mbx^\top \mbLambda(t)\nabla_{\tilde\mbx}f(\tilde\mbx)
\end{equation*}
for $f:\bbR^K\to\bbR$ twice continuously differentiable with derivatives having polynomial growth. Define the transformed source
\begin{equation}\label{app:eqn:g-definition}
    g(\tilde\mbx,t) = \nabla_{\tilde\mbx}\cdot\widetilde{\mbrho}(\tilde\mbx,t) - \tilde\mbx^\top\widetilde{\mbrho}(\tilde\mbx,t).
\end{equation}
A direct computation shows that
\begin{equation*}
    g(\tilde\mbx,t) = \eta(\mbm(t)+\mbS(t)^{1/2}\mbQ(t)\tilde\mbx,t),
\end{equation*}
where $\eta$ is the source term in \cref{app:eqn:helmholtz-stein-eqn}. Overall, in the $\tilde\mbx$ coordinates, the Poisson equation becomes
\begin{equation}\label{app:eqn:whitened-poisson-theorem}
    (\tilde{\cA}_t \chi(\cdot, t))(\tilde\mbx) = g(\tilde\mbx, t).
\end{equation}

By \Cref{app:assum:resid}, $\mbr$ is differentiable in $\mbx$, and both $\mbr$ and $D_{\mbx}\mbr$ have polynomial growth; the same therefore holds of $\widetilde{\mbrho}(\cdot,t)$ and $D_{\tilde\mbx}\widetilde{\mbrho}(\cdot,t)$. Stein's identity gives
\begin{equation*}
    \E_{\tilde\mbx \sim \cN(\mb0,\bbI_K)}\left[\nabla_{\tilde\mbx}\cdot \widetilde{\mbrho}(\tilde\mbx,t)\right] = \E_{\tilde\mbx \sim \cN(\mb0,\bbI_K)}\left[\tilde\mbx^\top \widetilde{\mbrho}(\tilde\mbx,t)\right],
\end{equation*}
so $g(\cdot, t)$ has mean zero under the OU invariant law.

In these whitened coordinates, the candidate from \cref{app:eqn:poisson-soln} takes the form
\begin{equation}\label{app:eqn:chi-representation-theorem}
    \chi(\tilde\mbx,t) = -\int_0^\infty (P_s^t g)(\tilde\mbx)\,ds, \qquad (P_s^t g)(\tilde\mbx) = \bbE\left[g\left(\tilde\mbz^{(t)}(s),t\right)\,\big|\,\tilde\mbz^{(t)}(0)=\tilde\mbx\right],
\end{equation}
where $(P_s^t)_{s \geq 0}$ is the OU semigroup associated with \cref{app:eqn:ou-process-theorem}.
Equivalently,
\begin{equation*}
    u(\mbx,t) = \chi\left( \mbQ(t)^\top\mbS(t)^{-1/2}(\mbx-\mbm(t)),t \right).
\end{equation*}

\medskip
\textbf{Step \ref{item:reg}}:
We first justify that \cref{app:eqn:chi-representation-theorem} is well-defined.
For all times $0 \le s < \infty$, the marginal distribution of the OU process
\eqref{app:eqn:ou-process-theorem} satisfies
\begin{equation*}
    \tilde\mbz^{(t)}(s) \overset{d}{=} \mbA_t(s)\tilde\mbx + \mbB_t(s)\mbxi, \qquad \mbxi \sim \cN(\mb0,\bbI_K),
\end{equation*}
where
\begin{equation*}
    \mbA_t(s) = e^{-s\mbLambda(t)/2}, \qquad \mbB_t(s) = \left(\bbI_K - e^{-s\mbLambda(t)}\right)^{1/2}.
\end{equation*}
Thus
\begin{equation*}
    (P_s^t g)(\tilde\mbx) = \bbE_{\mbxi \sim \cN(\mb0,\bbI_K)} \left[g\left(\mbA_t(s)\tilde\mbx+\mbB_t(s)\mbxi, t\right) \right].
\end{equation*}
Using $\E_{\mbxi \sim \cN(\mb0,\bbI_K)}[g(\mbxi,t)] = 0$, we may write
\begin{equation*}
    (P_s^t g)(\tilde\mbx) = \bbE_{\mbxi \sim \cN(\mb0,\bbI_K)} \left[ g\left(\mbA_t(s)\tilde\mbx+\mbB_t(s)\mbxi,t\right) - g(\mbxi,t) \right].
\end{equation*}
Hence,
\begin{equation*}
    \left|(P_s^t g)(\tilde\mbx)\right| \leq \bbE_{\mbxi \sim \cN(\mb0,\bbI_K)} \left| g\left(\mbA_t(s)\tilde\mbx+\mbB_t(s)\mbxi,t\right)-g(\mbxi,t) \right|.
\end{equation*}
Now decompose
\begin{equation*}
\begin{aligned}
    g\left(\mbA_t(s)\tilde\mbx+\mbB_t(s)\mbxi,t\right)-g(\mbxi,t) &= \Big( g\left(\mbA_t(s)\tilde\mbx+\mbB_t(s)\mbxi,t\right) - g\left(\mbB_t(s)\mbxi,t\right) \Big) \\
    &\quad + \Big( g\left(\mbB_t(s)\mbxi,t\right)-g(\mbxi,t) \Big).
\end{aligned}
\end{equation*}
Applying the mean value theorem to each term and using that $\nabla_{\tilde\mbx} g$ has at most polynomial growth,
there exist constants $C, m > 0$ such that
\begin{equation*}
    \|\nabla_{\tilde\mbx} g(\tilde\mbx,t)\|_2 \le C(1+\|\tilde\mbx\|_2^m).
\end{equation*}
Since $\mbxi$ is Gaussian, all of its moments are finite. Moreover,
\begin{equation*}
    \|\mbA_t(s)\|_{\mathrm{op}} \leq C_t e^{-c_t s}, \qquad \|\mbB_t(s)-\bbI_K\|_{\mathrm{op}} \leq C_t e^{-c_t s},
\end{equation*}
for some constants $C_t, c_t > 0$. It follows that
\begin{equation*}
    \left| \bbE\left[g\left(\tilde\mbz^{(t)}(s),t\right)\,\big|\,\tilde\mbz^{(t)}(0)=\tilde\mbx\right] \right| \leq C_t'(1+\|\tilde\mbx\|_2^{m+1})e^{-c_t s},
\end{equation*}
for some constant $C_t' > 0$. In particular, \cref{app:eqn:chi-representation-theorem} is finite for every $(\tilde\mbx, t)$.

We next show that, for every fixed $t$, $\chi(\cdot,t)$ is twice continuously differentiable in $\tilde\mbx$. We then establish joint continuity in $(\mbx,t)$ after transforming these derivatives back to the original coordinates. The estimates below also justify passage of the spatial derivatives, and hence of $\tilde{\cA}_t$, under the integral sign in step \ref{item:poisson}.

Using the notation introduced above, we can rewrite
\cref{app:eqn:chi-representation-theorem} as
\begin{equation}\label{app:eqn:chi-gaussian-representation-theorem}
    \chi(\tilde\mbx,t) = -\int_0^\infty \E_{\mbxi \sim \cN(\mb0,\bbI_K)} \left[ g\left(\mbA_t(s)\tilde\mbx+\mbB_t(s)\mbxi,t\right) \right] ds.
\end{equation}

We first study the gradient of $\chi$.
To do so, we must justify the exchange of the gradient and integral in \cref{app:eqn:chi-gaussian-representation-theorem}. Applying Stein's identity to the integrand of \cref{app:eqn:chi-gaussian-representation-theorem} yields 
\begin{equation*}
    F(\tilde\mbx,t,s) = \nabla_{\tilde\mbx}(P_s^t g)(\tilde\mbx) = \mbA_t(s)\mbB_t(s)^{-1} \bbE_{\mbxi \sim \cN(0,\bbI_K)} \left[ \mbxi \, g\left(\mbA_t(s)\tilde\mbx+\mbB_t(s)\mbxi,t\right) \right].
\end{equation*}
The polynomial growth of $g$ implies
\begin{equation*}
\begin{aligned}
    \|F(\tilde\mbx,t,s)\|_2 &\leq \|\mbA_t(s)\mbB_t(s)^{-1}\|_{\mathrm{op}} \bbE_{\mbxi \sim \cN(0,\bbI_K)} \left[ \|\mbxi\|_2\,|g\left(\mbA_t(s)\tilde\mbx+\mbB_t(s)\mbxi,t\right)| \right] \\
    &\leq C\|\mbA_t(s)\mbB_t(s)^{-1}\|_{\mathrm{op}} \left(1+\|\mbA_t(s)\tilde\mbx\|_2^m\right),
\end{aligned}
\end{equation*}
where we used that $\mbxi$ is Gaussian and $\mbB_t(s)$ is bounded.
Now, as $s\downarrow 0$, $\|\mbA_t(s)\mbB_t(s)^{-1}\|_{\mathrm{op}} = O(s^{-1/2})$,
while as $s\to\infty$,
\begin{equation*}
    \|\mbA_t(s)\|_{\mathrm{op}} \le C_t e^{-c_t s}, \quad \|\mbB_t(s)^{-1}\|_{\mathrm{op}} \le C_t.
\end{equation*}
Therefore
\begin{equation*}
    \|F(\tilde\mbx,t,s)\|_2 \leq \begin{cases} C(1+\|\tilde\mbx\|_2^m)s^{-1/2}, & 0 < s \leq 1,\\ C_t(1+\|\tilde\mbx\|_2^m)e^{-c_t s}, & s\geq 1, \end{cases}
\end{equation*}
and the right-hand side is integrable on $(0,\infty)$. Moreover, for every compact set $K \subset \bbR^K \times [0,T]$, this integrable bound can be chosen uniformly for $(\tilde\mbx,t)\in K$. By dominated convergence we deduce 
\begin{equation*}
    \nabla_{\tilde\mbx} \chi(\tilde\mbx,t) = -\int_0^\infty F(\tilde\mbx, t, s) ds.
\end{equation*}

We next study the Hessian of $\chi$.
Since $g$ is continuously differentiable in $\tilde\mbx$, for each fixed $s>0$ the map
$(\tilde\mbx,t)\mapsto F(\tilde\mbx,t,s)$ is continuously differentiable in $\tilde\mbx$, and
\begin{equation*}
    \nabla_{\tilde\mbx} F(\tilde\mbx,t,s) = \nabla_{\tilde\mbx}^2(P_s^t g)(\tilde\mbx) = \mbA_t(s)\mbB_t(s)^{-1} \bbE_{\mbxi \sim \cN(0,\bbI_K)} \left[ \mbxi \, \left(\nabla_{\tilde\mbx} g\left(\mbA_t(s)\tilde\mbx+\mbB_t(s)\mbxi,t\right)\right)^\top \right] \mbA_t(s).
\end{equation*}
Using that $\nabla_{\tilde\mbx}g$ has at most polynomial growth, we obtain
\begin{equation*}
\begin{aligned}
    \|\nabla_{\tilde\mbx} F(\tilde\mbx,t,s)\|_2&\leq\|\mbA_t(s)\mbB_t(s)^{-1}\|_{\mathrm{op}} \,\|\mbA_t(s)\|_{\mathrm{op}} \, \bbE_{\mbxi \sim \cN(0,\bbI_K)} \left[\|\mbxi\|_2\,\|\nabla_{\tilde\mbx} g\left(\mbA_t(s)\tilde\mbx+\mbB_t(s)\mbxi,t\right)\|_2 \right] \\
    &\leq C\|\mbA_t(s)\mbB_t(s)^{-1}\|_{\mathrm{op}}\,\|\mbA_t(s)\|_{\mathrm{op}} \left(1+\|\mbA_t(s)\tilde\mbx\|_2^m\right).
\end{aligned}
\end{equation*}
Hence,
\begin{equation*}
    \|\nabla_{\tilde\mbx} F(\tilde\mbx,t,s)\|_2 \leq \begin{cases} C(1+\|\tilde\mbx\|_2^m)s^{-1/2}, & 0 < s \leq 1,\\ C_t(1+\|\tilde\mbx\|_2^m)e^{-c_t s}, & s\ge 1, \end{cases}
\end{equation*}
which is again integrable on $(0,\infty)$. Again, for every compact set $K \subset \bbR^K \times [0,T]$, this integrable bound may be chosen uniformly for $(\tilde\mbx,t)\in K$, and dominated convergence yields
\begin{equation*}
    \nabla_{\tilde\mbx}^2 \chi(\tilde\mbx,t) = -\int_0^\infty \nabla_{\tilde\mbx}F(\tilde\mbx,t,s)\,ds.
\end{equation*}
Hence, $\chi(\cdot,t)$ is twice continuously differentiable in $\tilde\mbx$ for every $t$.

It remains to establish continuity after transforming back to the original coordinates. A continuous choice of the eigenvector matrix $\mbQ(t)$ is not required. Indeed,
\begin{equation*}
    \mbQ(t)\mbA_t(s)\mbQ(t)^\top = e^{-s\mbC(t)/2}, \quad  \mbQ(t)\mbB_t(s)\mbQ(t)^\top = \left(\bbI_K-e^{-s\mbC(t)}\right)^{1/2}.
\end{equation*}
Consequently, after undoing the rotation, the expressions for the derivatives of $u$ depend on $t$ through $\mbm(t)$, $\mbS(t)$, $\mbr$, and the continuous matrix functions of $\mbC(t)$ displayed above, rather than through the particular choice of $\mbQ(t)$.

Continuity and positive definiteness of $\mbS(t)$ and $\mbSigma(t)$ imply uniform upper and lower spectral bounds for $\mbC(t)$ on $[0,T]$. Therefore the bounds above may be chosen locally uniformly in $(\mbx,t)$. Dominated convergence in the Gaussian expectation and in $s$ then gives joint continuity of the first and second spatial derivatives in $\mbx$ coordinates.

\medskip
\textbf{Step \ref{item:poisson}}:
We now verify that $u$ satisfies the Poisson equation \cref{app:eqn:weighted-poisson-theorem} pointwise.
By step \ref{item:whiten}, it suffices to show that $\chi$ satisfies the Poisson equation in the transformed coordinates \eqref{app:eqn:whitened-poisson-theorem}.

The estimates established in step \ref{item:reg} show that $(P_s^t g)(\tilde\mbx)$ is twice continuously differentiable in $\tilde\mbx$.
For each $s>0$, $(P_s^t g)$ is also continuously differentiable in $s$ and satisfies the backward Kolmogorov equation
\begin{equation*}
    \frac{\partial}{\partial s}(P_s^t g)(\tilde\mbx) = (\tilde{\cA}_t P_s^t g)(\tilde\mbx), \qquad s>0.
\end{equation*}

Moreover, the spatial derivatives of $(P_s^t g)$ admit integrable bounds in $s$, locally uniformly in $\tilde\mbx$ for fixed $t$.
Therefore the operator $\tilde{\cA}_t$ may be passed under the integral sign in \cref{app:eqn:chi-representation-theorem}, yielding
\begin{equation*}
\begin{aligned}
    (\tilde{\cA}_t \chi(\cdot,t))(\tilde\mbx) &= -\int_0^\infty (\tilde{\cA}_t P_s^t g)(\tilde\mbx)ds\\
    &= -\int_0^\infty \frac{\partial}{\partial s}(P_s^t g)(\tilde\mbx)ds\\
    &= g(\tilde\mbx,t)-\lim_{R\to\infty}(P_R^t g)(\tilde\mbx).
\end{aligned}
\end{equation*}
By the tail estimate established in step \ref{item:reg},
\begin{equation*}
    |(P_R^t g)(\tilde\mbx)| \leq C_t'(1+\|\tilde\mbx\|_2^{m+1})e^{-c_tR} \to 0 \qquad\text{as }R\to\infty,
\end{equation*}
for each $(\tilde\mbx,t)$.
Hence $(\tilde{\cA}_t\chi(\cdot,t))(\tilde\mbx)=g(\tilde\mbx,t)$ for every $(\tilde\mbx,t)$, and consequently
\begin{equation*}
    u(\mbx,t) = \chi\left( \mbQ(t)^\top\mbS(t)^{-1/2}(\mbx-\mbm(t)),t \right)
\end{equation*}
satisfies \cref{app:eqn:weighted-poisson-theorem} pointwise. The gradient component of the Helmholtz decomposition is then $\nabla_{\mbx}\psi(\mbx,t)=\tfrac{1}{2}\nabla_{\mbx}u(\mbx,t)$.

\medskip
\textbf{Step \ref{item:unique}}:
It remains to identify the constructed vector fields with the components of the weak Helmholtz decomposition in
\Cref{app:thm:hilbert-projection}. 
The estimates in step \ref{item:reg} show that $\psi$, $\nabla_{\mbx}\psi$, and $\mbh$ have polynomial growth. 
Since $q(\mbx,t)$ is Gaussian, $\mbSigma(t)\nabla_{\mbx}\psi(\mbx,t)$ and $\mbh(\mbx,t)$ belong to $\cL^2(q(\mbx, t), \bbR^K, \mbSigma(t))^K$.

First, by appealing to integration by parts, the pointwise identity $\nabla_{\mbx}\cdot \left[q(\mbx,t)\mbh(\mbx,t)\right] = 0$ implies that $\mbh(\mbx,t)$ is $q$-divergence-free in the weak sense \cref{app:eqn:div-free-weak}. 
We next show that $\mbSigma(t)\nabla_{\mbx}\psi(\cdot,t)$ is orthogonal to all weakly $q$-divergence-free vector fields.

Fix $\mbv \in \cV_t$, and define an infinitely differentiable function $\zeta_R$ satisfying $\zeta_R=1$ on the ball of radius $R$, $\zeta_R=0$ outside the ball of radius $3R/2$, and $\|\nabla_{\mbx}\zeta_R\|_2 \leq \frac{C}{R}$.
$\zeta_R(\mbx)\psi(\mbx,t)$ is twice continuously differentiable and compactly supported on the ball of radius $3R/2$. 
We may therefore choose $\{\varphi_{R,k}\}_{k=1}^{\infty}$ infinitely differentiable and compactly supported on the ball of radius $2R$ such that
\begin{equation*}
    \nabla_{\mbx}\varphi_{R,k} \to \nabla_{\mbx}\left[    \zeta_R\psi(\cdot,t) \right]
\end{equation*}
uniformly. Since $\mbv\in\cV_t$,
\begin{equation*}
    \int
    \mbv(\mbx)^\top\nabla_{\mbx}\varphi_{R,k}(\mbx)
    q(\mbx,t)d\mbx
    =
    0
\end{equation*}
for every $k$. 
Cauchy--Schwarz allows us to pass to the limit as $k\to\infty$, producing
\begin{equation*}
    \int \mbv(\mbx)^\top \nabla_{\mbx}\left[   \zeta_R(\mbx)\psi(\mbx,t) \right] q(\mbx,t)d\mbx = 0.
\end{equation*}
The polynomial growth of $\psi$ and $\nabla_{\mbx}\psi$, together with the Gaussian tails of $q(\mbx,t)$, gives
\begin{equation*}
\left\| \mbSigma(t) \nabla_{\mbx}\left[ \zeta_R(\mbx)\psi(\mbx,t) \right] - \mbSigma(t)\nabla_{\mbx}\psi(\mbx,t) \right\|_{q, t} \to 0
\end{equation*}
as $R\to\infty$.
Taking the limit in the preceding identity therefore gives
$\left\langle \mbv(\mbx), \mbSigma(t)\nabla_{\mbx}\psi(\mbx,t) \right\rangle_{q, t} = 0$. 
Since this holds for every $\mbv \in \cV_t$, we have the desired result.

Thus,
\begin{equation*}
\mbr(\mbx,t) = \mbSigma(t)\nabla_{\mbx}\psi(\mbx,t) + \mbh(\mbx,t)
\end{equation*}
is precisely the unique orthogonal decomposition supplied by \Cref{app:thm:hilbert-projection}. 
In particular, the two vector fields are unique up to $q(\mbx,t)$-almost-everywhere equality among all square-integrable weak Helmholtz decompositions. 
Because $q(\mbx,t) > 0$ for every $\mbx$ and the constructed vector fields are continuous, almost-everywhere uniqueness implies pointwise uniqueness.
\end{proof}

\Cref{app:thm:helm-gaussian} implies that for Gaussian marginals and a state-independent diffusion coefficient, the Helmholtz decomposition exists and admits the probabilistic solution
\begin{equation}\label{app:eq:prob-rep}
\begin{aligned}
    \nabla_{\mbx}\psi(\mbx,t)
    &=
    \frac{1}{2}\mbS(t)^{-1/2}\mbQ(t)
    \nabla_{\tilde\mbx}\chi\left(
        \mbQ(t)^\top\mbS(t)^{-1/2}(\mbx-\mbm(t)),t
    \right),\\
    \nabla_{\tilde\mbx}\chi(\tilde\mbx,t)
    &=
    -\int_0^\infty
    \mbA_t(s)\mbB_t(s)^{-1}
    \bbE_{\mbxi\sim\cN(\mb0,\bbI_K)}
    \left[
        \mbxi\,
        g\left(
            \mbA_t(s)\tilde\mbx+
            \mbB_t(s)\mbxi,t
        \right)
    \right]ds,\\
    \mbA_t(s)
    &=
    e^{-s\mbLambda(t)/2},
    \qquad
    \mbB_t(s)
    =
    \left(\bbI_K-e^{-s\mbLambda(t)}\right)^{1/2},\\
    \mbC(t)
    &=
    \mbS(t)^{-1/2}\mbSigma(t)\mbS(t)^{-1/2}
    =
    \mbQ(t)\mbLambda(t)\mbQ(t)^\top,\\
    g(\tilde\mbx,t)
    &=
    \nabla_{\tilde\mbx}\cdot
    \widetilde{\mbrho}(\tilde\mbx,t)
    -
    \tilde\mbx^\top\widetilde{\mbrho}(\tilde\mbx,t),\\
    \widetilde{\mbrho}(\tilde\mbx,t)
    &=
    \mbQ(t)^\top\mbS(t)^{-1/2}
    \mbr\left(
        \mbm(t)+\mbS(t)^{1/2}\mbQ(t)\tilde\mbx,t
    \right).
\end{aligned}
\end{equation}

\subsection{Polynomial approximation of the optimal correction}
\label{app:subsec:resid-approx}

In \Cref{subsec:helmholtz-decomp} we showed that at fixed marginals $\{q(\mbx, t)\}_{0 \leq t \leq T}$, the $q$-divergence-free component $\mbh$ of the weighted Helmholtz decomposition is the vector field $\mbv$ that minimizes
\begin{equation}\label{app:eq:fixed-marginal-objective}
\int_0^T \bbE_{\bbQ} \left\|\mbSigma(t)^{-1/2} \{\mbr(\mbx(t), t)-\mbv(\mbx(t), t)\}\right\|^2dt.
\end{equation}
However, even when $q(\mbx, t)$ is Gaussian, the correction generally admits only the probabilistic solution \cref{app:eqn:helmholtz-soln}.
In our preliminary experiments, we found that estimating \cref{app:eqn:helmholtz-soln} directly has prohibitively high variance.
We make the problem tractable by restricting the divergence-free correction $\mbv$ to a finite-dimensional subspace comprised of polynomials of degree at most $\ell$.
Within this subspace, the minimizer of \cref{app:eq:fixed-marginal-objective} can be computed coefficient by coefficient in a Hermite polynomial basis.

Throughout this subsection, we assume only that $(\mbm(t),\mbS(t))_{0 \leq t \leq T}$, $\mbr(\mbx, t)$, $\mbSigma(t)$ satisfy the regularity conditions in \Cref{app:assum:resid}.

\paragraph{Whiten and rotate.}
First, we define the whitened coordinates $\tilde\mbx$ just as we did in the proof of \Cref{app:thm:helm-gaussian}:
\begin{equation}\label{app:eq:hermite-coordinates}
\tilde\mbx = \mbQ(t)^\top\mbS(t)^{-1/2}(\mbx-\mbm(t)), \quad \widetilde{\mbrho}(\tilde\mbx, t) =\mbQ(t)^\top\mbS(t)^{-1/2}\mbr(\mbm(t)+\mbS(t)^{1/2}\mbQ(t)\tilde\mbx, t),
\end{equation}
where $\mbQ(t)\mbLambda(t)\mbQ(t)^\top$ is the eigendecomposition of $\mbC(t)=\mbS(t)^{-1/2}\mbSigma(t)\mbS(t)^{-1/2}$.
In the whitened coordinates, $\tilde\mbx\sim\cN(\mb0,\bbI_K)$.

Suppose $\mbnu(\tilde\mbx,t)$ is divergence-free with respect to the standard Gaussian density.
In the original coordinates, this vector field is
\begin{equation}\label{app:eq:hermite-unwhiten}
\mbv(\mbx, t) =\mbS(t)^{1/2}\mbQ(t)\mbnu\left(\mbQ(t)^\top\mbS(t)^{-1/2}(\mbx-\mbm(t)), t\right).
\end{equation}
The objective and constraint become
\begin{align}
&\bbE_{\tilde\mbx\sim\cN(\mb0,\bbI_K)}\left[(\widetilde{\mbrho}(\tilde\mbx,t)-\mbnu(\tilde\mbx,t))^\top\mbLambda(t)^{-1}(\widetilde{\mbrho}(\tilde\mbx,t)-\mbnu(\tilde\mbx,t))\right]\label{app:eq:hermite-objective}\\
&\nabla_{\tilde\mbx}\cdot\mbnu(\tilde\mbx,t)-\tilde\mbx^\top\mbnu(\tilde\mbx,t)=0.\label{app:eq:hermite-div-free}
\end{align}
Indeed, \cref{app:eq:hermite-div-free} is exactly $\nabla_{\mbx}\cdot[q(\mbx,t)\mbv(\mbx,t)]=0$ after the change of variables.

\paragraph{Closed-form Hermite projection.}
Let $\mathrm{He}_d(u)$, $d\in\bbN$, be the $d$th unnormalized probabilists' Hermite polynomial.
The collection $\{\mathrm{He}_d(u)/\sqrt{d!}\}_{d=0}^{\infty}$ constitutes an orthonormal basis of the Hilbert space $\cL^2(\cN(0,1);\bbR)$, which is comprised of functions $f:\bbR\to\bbR$ satisfying $\bbE_{u\sim\cN(0,1)}[f(u)^2]<\infty$ with the inner product $\langle f,g\rangle=\bbE_{u\sim\cN(0,1)}[f(u)g(u)]$.

We can extend the univariate Hermite polynomial basis to an orthonormal basis of $\cL^2(\cN(\mb0,\bbI_K);\bbR)$, comprised of functions $f:\bbR^K\to\bbR$ satisfying $\bbE_{\tilde\mbx\sim\cN(\mb0,\bbI_K)}[f(\tilde\mbx)^2]<\infty$ with the analogous inner product.
Define the collection of multi-indices
\begin{equation*}
\cN_{K, \ell} = \left\{ \mbalpha \in \{0, 1, \ldots, \ell\}^K \bigg| \sum_{j=1}^K \mbalpha_j \leq \ell\right\}, \quad \cN_{K, +\infty} = \{0, 1, \ldots\}^K, \quad \ell \geq 0.
\end{equation*}
For a multi-index $\mbalpha \in \cN_{K, +\infty}$, we define the corresponding normalized multivariate probabilists' Hermite polynomial
\begin{equation}\label{app:eq:hermite-basis}
\hat{H}_{\mbalpha}(\tilde\mbx) =\prod_{i=1}^K \frac{\mathrm{He}_{\mbalpha_i}(\tilde\mbx_i)}{\sqrt{\mbalpha_i!}}.
\end{equation}
$\{\hat{H}_{\mbalpha}(\tilde\mbx)\}_{\mbalpha\in\cN_{K,+\infty}}$ forms an orthonormal basis of $\cL^2(\cN(\mb0,\bbI_K);\bbR)$.

By \Cref{app:assum:resid}, $\widetilde{\mbrho}_i(\tilde\mbx,t)$, $1\leq i\leq K$, has polynomial growth and hence is in $\cL^2(\cN(\mb0,\bbI_K);\bbR)$.
Therefore, we can expand the rotated residual in the multivariate Hermite basis as
\begin{equation}\label{app:eq:hermite-coefficients}
\widetilde{\mbrho}_i(\tilde\mbx,t)=\sum_{\mbalpha\in\cN_{K,+\infty}}\widetilde{\mbb}_{\mbalpha,i}(t)\hat{H}_{\mbalpha}(\tilde\mbx), \qquad \widetilde{\mbb}_{\mbalpha,i}(t)=\bbE_{\tilde\mbx\sim\cN(\mb0,\bbI_K)}\left[\widetilde{\mbrho}_i(\tilde\mbx,t)\hat{H}_{\mbalpha}(\tilde\mbx)\right].
\end{equation}
We choose this basis expansion because of the identity
\begin{equation}\label{app:eq:hermite-raising}
(\partial_{\tilde\mbx_i}-\tilde\mbx_i)\hat{H}_{\mbalpha}(\tilde\mbx) =-\sqrt{\mbalpha_i+1}\,\hat{H}_{\mbalpha+\mbe_i}(\tilde\mbx).
\end{equation}
The Gaussian weighted divergence raises the Hermite degree by one and creates no lower-order terms.
Rather than keeping the infinitely many basis vectors in the expansion of $\widetilde{\mbrho}_i(\tilde\mbx, t)$, we truncate the expansion \cref{app:eq:hermite-coefficients} to $\{\hat{H}_{\mbalpha}(\tilde\mbx)\}_{\mbalpha \in \cN_{K, \ell}}$, which span the $\binom{K + \ell}{\ell}$-dimensional subspace of polynomials having degree at most $\ell$.
In \Cref{app:prop:hermite-projection}, we show that by doing so, we can obtain the optimal correction to the reference posterior drift restricted to this subspace.
\begin{proposition}[Hermite projection]\label{app:prop:hermite-projection}
Among all polynomial fields $\mbnu(\tilde\mbx, t)$ of degree at most $\ell$ satisfying \cref{app:eq:hermite-div-free}, the unique minimizer of \cref{app:eq:hermite-objective} is
\begin{equation}\label{app:eq:hermite-solution}
\begin{aligned}
\mu_{\mbbeta}(t)&=\frac{\sum_{i=1}^K\sqrt{\mbbeta_i}\, \widetilde{\mbb}_{\mbbeta-\mbe_i,i}(t)} {\sum_{i=1}^K\lambda_i(t)\mbbeta_i}, &&\mbbeta \in \cN_{K, \ell+1} \setminus \cN_{K, 0},\\
\widetilde{\mbc}_{\mbalpha,i}(t) &=\widetilde{\mbb}_{\mbalpha,i}(t) -\lambda_i(t)\sqrt{\mbalpha_i+1}\, \mu_{\mbalpha+\mbe_i}(t), &&\mbalpha \in \cN_{K, \ell},\\
\mbnu^{(\ell)}(\tilde\mbx, t) &=\sum_{\mbalpha \in \cN_{K, \ell}} \widetilde{\mbc}_{\mbalpha}(t)\hat{H}_{\mbalpha}(\tilde\mbx),
\end{aligned}
\end{equation}
where a coefficient with a negative multi-index is zero.
The optimal correction $\mbh^{(\ell)}(\mbx, t)$ in the original coordinates is obtained from $\mbnu^{(\ell)}(\tilde\mbx, t)$ using \cref{app:eq:hermite-unwhiten}.
\end{proposition}

\begin{proof}
Note that since $\mbnu$ is a polynomial field, its components are in $\cL^2(\cN(\mb0,\bbI_K);\bbR)$.
Expanding $\mbnu$ in the multivariate Hermite basis $\mbnu_i(\tilde\mbx, t)=\sum_{|\mbalpha|\leq\ell}\widetilde{\mbc}_{\mbalpha,i}(t)\hat{H}_{\mbalpha}(\tilde\mbx)$ and applying \cref{app:eq:hermite-raising} to $\mbnu(\tilde\mbx, t)$ gives
\begin{equation}\label{app:eq:hermite-stein-action}
\nabla_{\tilde\mbx}\cdot\mbnu(\tilde\mbx, t) -\tilde\mbx^\top\mbnu(\tilde\mbx, t) =-\sum_{1 \leq |\mbbeta| \leq \ell+1} \left( \sum_{i=1}^K\sqrt{\mbbeta_i}\, \widetilde{\mbc}_{\mbbeta-\mbe_i,i}(t) \right)\hat{H}_{\mbbeta}(\tilde\mbx).
\end{equation}
Thus divergence freedom gives one scalar constraint for each $\mbbeta$.
Moreover, every coefficient $\widetilde{\mbc}_{\mbalpha,i}(t)$ occurs in exactly one constraint, namely $\mbbeta=\mbalpha+\mbe_i$.

By Parseval's identity, the part of the objective that depends on the correction is
\begin{equation*}
\sum_{|\mbalpha|\leq\ell} (\widetilde{\mbb}_{\mbalpha}(t)- \widetilde{\mbc}_{\mbalpha}(t))^\top \mbLambda(t)^{-1} (\widetilde{\mbb}_{\mbalpha}(t)- \widetilde{\mbc}_{\mbalpha}(t)).
\end{equation*}
For a fixed $\mbbeta$, the relevant coefficients solve
\begin{equation*}
\min_{\{c_i\}}\frac{1}{2}\sum_{i:\mbbeta_i>0}\frac{1}{\lambda_i(t)}\left(\widetilde{\mbb}_{\mbbeta-\mbe_i,i}(t)-c_i\right)^2 \quad\text{subject to}\quad \sum_{i:\mbbeta_i>0}\sqrt{\mbbeta_i}\,c_i=0.
\end{equation*}
Introducing the Lagrange multiplier $\mu_{\mbbeta}(t)$ and differentiating the Lagrangian yields the first-order conditions
\begin{equation*}
c_i = \widetilde{\mbb}_{\mbbeta-\mbe_i,i}(t) -\lambda_i(t)\sqrt{\mbbeta_i}\,\mu_{\mbbeta}(t).
\end{equation*}
Substituting this expression into the constraint yields
\begin{equation*}
\begin{aligned}
&\sum_i\sqrt{\mbbeta_i} c_i = \sum_i\sqrt{\mbbeta_i} \widetilde{\mbb}_{\mbbeta-\mbe_i,i}(t) - \mu_{\mbbeta}(t)\sum_i\lambda_i(t)\mbbeta_i=0,\\
\implies& \mu_{\mbbeta}(t) =\frac{\sum_{i=1}^K\sqrt{\mbbeta_i}\widetilde{\mbb}_{\mbbeta-\mbe_i,i}(t)} {\sum_{i=1}^K\lambda_i(t)\mbbeta_i}.
\end{aligned}
\end{equation*}
Because every $\lambda_i(t)>0$, the objective is strictly convex.
Consequently, this feasible stationary point is the unique minimizer.
\end{proof}

Recall from \Cref{app:subsec:helmholtz-compute} that the weighted Helmholtz decomposition can be obtained by solving a Poisson equation \cref{app:eqn:poisson} defined by the infinitesimal generator of the Langevin diffusion \cref{app:eqn:langevin-diffusion}.
In the whitened coordinates, the Langevin diffusion is a time-rescaled Ornstein--Uhlenbeck process \cref{app:eqn:ou-process-theorem}, and the eigenfunctions of the  infinitesimal generator are the Hermite polynomials.
Indeed, evaluating \cref{app:eq:hermite-raising} at $\mbalpha - \mbe_i$ and combining it with $\partial_{\mbx_i} \hat{H}_{\mbalpha}(\tilde\mbx) = \sqrt{\mbalpha_i} \hat{H}_{\mbalpha - \mbe_i}(\tilde\mbx)$ produces
\begin{equation*}
    (\partial_{\tilde\mbx_i}-\tilde\mbx_i)\partial_{\mbx_i} \hat{H}_{\mbalpha}(\tilde\mbx) = \sqrt{\mbalpha_i} (\partial_{\tilde\mbx_i}-\tilde\mbx_i) \hat{H}_{\mbalpha - \mbe_i}(\tilde\mbx) = - \mbalpha_i \hat{H}_{\mbalpha}(\tilde\mbx).
\end{equation*}
Hence, \Cref{app:prop:hermite-projection} can be understood as expanding the left and right-hand sides of the Poisson equation in the orthonormal basis of eigenfunctions and then solving for the coefficients belonging to $\cN_{K, \ell+1}$ (indexed by $\mbbeta$).
We then differentiate the result to obtain the gradient field $\nabla_{\tilde\mbx} \chi^{(\ell)}(\tilde\mbx,t)$.

\paragraph{Two useful special cases.}
When $\ell = 0$, $\mbbeta=\mbe_i$, \cref{app:eq:hermite-solution} gives $\widetilde{\mbc}_{\mb0,i}(t)=0$, so a compatible correction has no constant term and cannot change the prescribed mean path.
For $\ell=1$, write $\mbnu^{(1)}(\tilde\mbx, t)=\widetilde{\mbA}(t)\tilde\mbx$ and set
\begin{equation*}
\widetilde{\mbB}_{ij}(t) =\widetilde{\mbb}_{\mbe_j,i}(t) =\bbE_{\tilde\mbx\sim\cN(\mb0,\bbI_K)}\left[\widetilde{\mbrho}_i(\tilde\mbx,t)\tilde\mbx_j\right].
\end{equation*}
By Stein's identity, $\widetilde{\mbB}(t)=\bbE_{\tilde\mbx\sim\cN(\mb0,\bbI_K)}[D_{\tilde\mbx}\widetilde{\mbrho}(\tilde\mbx,t)]$.
The solution reduces to
\begin{equation}\label{app:eq:hermite-degree1}
\widetilde{\mbA}_{ij}(t) = \frac{\lambda_j(t)\widetilde{\mbB}_{ij}(t) -\lambda_i(t)\widetilde{\mbB}_{ji}(t)}{\lambda_i(t)+\lambda_j(t)}, \qquad \widetilde{\mbA}_{ii}(t)=0.
\end{equation}
Hence the linear correction is skew-symmetric in whitened coordinates.
If $\mbLambda(t)$ is proportional to the identity, it is simply the antisymmetric part of the averaged Jacobian.

\paragraph{Monte Carlo.}
Computing the polynomial approximation to the optimal divergence-free vector field requires computing the coefficients of $\widetilde{\mbrho}_i(\tilde\mbx, t)$ in the Hermite basis.
Generally, these coefficients cannot be computed in closed-form.
However, they do have the unbiased estimator
\begin{equation}\label{app:eq:hermite-mc}
\widehat{\widetilde{\mbb}}_{\mbalpha,i}(t)=\frac{1}{N_{\mathrm{mc}}}\sum_{n=1}^{N_{\mathrm{mc}}}\widetilde{\mbrho}_i(\tilde\mbx^{(n)},t)\hat{H}_{\mbalpha}(\tilde\mbx^{(n)}), \qquad \tilde\mbx^{(n)}\overset{\mathrm{i.i.d.}}{\sim}\cN(\mb0,\bbI_K).
\end{equation}
Substituting any estimated coefficients into \cref{app:eq:hermite-solution} still makes \cref{app:eq:hermite-stein-action} vanish identically.
Thus, the estimated field is $q$-divergence-free, although it will generally not be the optimal degree $\ell$ polynomial $q$-divergence-free vector field.
In the stochastic VI algorithm, the batch of points $\mbx_{j, k} \overset{\mathrm{i.i.d.}}{\sim} q(\mbx, t_j)$ used to evaluate the path KL $\cL_{\mathrm{path}}$ and the batch used to compute the correction to the reference drift at these points should be independent.

\paragraph{Gauss--Hermite quadrature.}
Alternatively, for low dimensional latent spaces, we can use Gauss--Hermite quadrature to estimate the Hermite coefficients of the residual vector field. If $N_{\mathrm{quad}}$ is the number of quadrature points along each latent dimension, the estimator is
\begin{equation}\label{app:eq:hermite-quad}
\widehat{\widetilde{\mbb}}_{\mbalpha,i}(t)=\frac{1}{N_{\mathrm{quad}}^K}\sum_{n=1}^{N_{\mathrm{quad}}^K} w_n \widetilde{\mbrho}_i(\tilde\mbx_n, t)\hat{H}_{\mbalpha}(\tilde\mbx_n)
\end{equation}
where the sum is taken over the Cartesian product of one-dimensional quadrature points $\tilde\mbx_n$, and $w_n$ is the product of the corresponding one-dimensional quadrature weights.
When $\widetilde{\mbrho}_i(\tilde\mbx, t) \hat{H}_{\mbalpha}(\tilde\mbx)$ is a polynomial of degree at most $2 N_{\mathrm{quad}} - 1$, Gauss--Hermite quadrature computes the Hermite coefficient $\widetilde{\mbb}_{\mbalpha,i}(t)$ exactly.


\paragraph{Stein estimand.}
If derivatives of the residual are available, Stein's identity provides an alternate estimator for $\widetilde{\mbb}_{\mbalpha,i}(t)$.
For any multi-index $\mbkappa$ such that $\mbkappa_j \leq \mbalpha_j, 1 \leq j \leq K$, repeated integration by parts gives
\begin{equation}\label{app:eq:stein-estimator}
\widetilde{\mbb}_{\mbalpha,i}(t)=\sqrt{\frac{(\mbalpha-\mbkappa)!}{\mbalpha!}}\,\bbE_{\tilde\mbx\sim\cN(\mb0,\bbI_K)}\left[\partial_{\tilde\mbx}^{\mbkappa}\widetilde{\mbrho}_i(\tilde\mbx,t)\hat{H}_{\mbalpha-\mbkappa}(\tilde\mbx)\right].
\end{equation}
The expectation in \cref{app:eq:stein-estimator} can be approximated using either Monte Carlo or Gauss--Hermite quadrature.
The choice $\mbkappa=\mb0$ recovers \cref{app:eq:hermite-mc}.
If derivatives through order $|\mbalpha|$ are available, choosing $\mbkappa=\mbalpha$ removes the Hermite polynomial entirely:
\begin{equation*}
\widetilde{\mbb}_{\mbalpha,i}(t) =\frac{1}{\sqrt{\mbalpha!}} \bbE_{\tilde\mbx\sim\cN(\mb0,\bbI_K)}\left[\partial_{\tilde\mbx}^{\mbalpha} \widetilde{\mbrho}_i(\tilde\mbx,t) \right].
\end{equation*}
Under the twice-differentiability condition in \Cref{app:assum:resid}, one may use any $\mbkappa \in \cN_{K, 2}$.
The derivatives with respect to $\tilde\mbx$ can be obtained from spatial derivatives of $\mbr$ through the affine change of coordinates in \cref{app:eq:hermite-coordinates} and evaluated using automatic differentiation.

Although the normalized Hermite polynomials have unit second moments, their higher moments grow rapidly with their degree.
Integration by parts can therefore reduce the variance of the coefficient estimates by replacing a high-degree Hermite factor with derivatives of the residual.
This reduction is not guaranteed, since the derivatives may themselves have high variance.
The choice of $\mbkappa$ trades the degree of the remaining Hermite factor against the cost and variability of higher-order derivatives.
Regardless of this choice, substituting the resulting coefficient estimates into \cref{app:eq:hermite-solution} preserves the divergence-free constraint exactly.

\paragraph{Computational complexity.}
There are $K\binom{K+\ell}{\ell}$ vector coefficients.
Once they have been estimated, \cref{app:eq:hermite-solution} touches each coefficient a constant number of times and costs $O(K\binom{K+\ell}{\ell})$, rather than $O(K \binom{K+\ell}{\ell}^3)$ from a dense linear solve.
Including residual evaluations, Hermite features, and eigendecomposition, the cost at one time point is
\begin{equation*}
O\left(N_{\mathrm{est}}C_{\mbr} + N_{\mathrm{est}} K\binom{K+\ell}{\ell}+K^3\right),
\end{equation*}
where $C_{\mbr}$ is the cost of one residual evaluation and $N_{\mathrm{est}}$ is the number of Monte Carlo samples or quadrature points used to approximate the Hermite coefficients.
When $\ell$ is fixed, $\binom{K+\ell}{\ell} = \Theta(K^{\ell})$, so the coefficient-dependent term scales as $K^{\ell+1}$.

The Stein estimand \cref{app:eq:stein-estimator} additionally requires derivatives of the transformed residual.
For the full-order choice $\mbkappa=\mbalpha$, the highest-order derivative tensor has $K^\ell$ entries for each of the $K$ output components, and hence $K^{\ell+1}$ entries in total.
Using automatic differentiation, computing these derivatives costs $O(K^{\ell+1} C_{\mbr})$.

\subsection{Verification of assumptions for the Helmholtz-SDE variational family}
\label{app:subsec:verify-assumptions}

In \Cref{app:subsec:helmholtz-gaussian,app:subsec:resid-approx} we discussed how to obtain a correction $\mbh^{(\ell)}$ to the reference posterior drift from the weighted Helmholtz decomposition of the residual vector field $\mbr = \mbf_p - \mbf_q$.
It remains to establish that this correction produces a valid variational SDE.
In this subsection, we verify that the Gaussian variational family used by Helmholtz-SDE satisfies the assumptions in \Cref{app:subsec:assumptions}.

Let $q(\mbx,t)=\cN(\mbm(t),\mbS(t))$.
We assume only that $(\mbm(t),\mbS(t))_{0 \leq t \leq T}$, $\mbr(\mbx, t)$, $\mbSigma(t)$ satisfy the regularity conditions in \Cref{app:assum:resid}. \Cref{app:assum:regularity,app:assum:densities,app:ass:path-ac,app:ass:finite-kl,app:ass:fpk-unique} are verified directly from \Cref{app:assum:resid}.

Helmholtz-SDE uses the reference drift
\begin{equation*}
    \mbf_q(\mbx,t)  = \frac{d}{dt}\mbm(t) + \left\{ \left(\frac{d}{dt}\mbL(t)\right)\mbL(t)^{-1} -\frac{1}{2}\mbSigma(t)\mbS(t)^{-1} \right\} (\mbx-\mbm(t)),
\end{equation*}
where $\mbL(t)$ is the symmetric square root of $\mbS(t)$.
The full posterior drift is $\mbf_q(\mbx,t)+\mbh^{(\ell)}(\mbx,t)$, where $\mbh^{(\ell)}$ is the degree $\ell$ correction constructed in \Cref{app:subsec:resid-approx}.  
The posterior initial density is $q(\mbx,0)$, and the prior initial distribution is $\cN(\mbmu_0,\mbV_0)$, with $\mbV_0$ positive definite.
The verification below applies for $0 \leq \ell \leq 3$.

Because $[0,T]$ is compact, the eigenvalues of $\mbS(t)$ and $\mbSigma(t)$ are bounded above and bounded away from zero.  
The Hermite coefficients defining $\mbh^{(\ell)}$ are continuous, and hence bounded, in $t$.  
Indeed, continuity follows from \Cref{app:assum:resid}, dominated convergence under the Gaussian marginals, and the finite-dimensional projection in \Cref{app:prop:hermite-projection}.

\paragraph{Existence and uniqueness.}
Introduce the whitened variable $\mbz$
\begin{equation*}
    \mbz=\mbL(t)^{-1}(\mbx-\mbm(t)), \qquad \mbnu^{(\ell)}(\mbz,t) =\mbL(t)^{-1} \mbh^{(\ell)}(\mbm(t)+\mbL(t)\mbz,t).
\end{equation*}
For the symmetric-square-root reference drift, the terms involving $d\mbL(t)/dt$ cancel under this change of variables. 
Consequently, the posterior SDE becomes
\begin{equation}
\label{app:eq:whitened-posterior-sde}
\begin{aligned}
    d\mbz(t) &= \left\{ -\frac{1}{2}\mbC(t)\mbz(t) +\mbnu^{(\ell)}(\mbz(t),t) \right\}dt+\mbL(t)^{-1}\mbSigma(t)^{1/2}d\mbw(t),\\\mbC(t) &=\mbS(t)^{-1/2}\mbSigma(t)\mbS(t)^{-1/2}.
\end{aligned}
\end{equation}
The drift in \cref{app:eq:whitened-posterior-sde} is polynomial in $\mbz$, and is therefore locally Lipschitz.  
This gives a pathwise-unique strong solution up to its explosion time.

It remains to rule out explosion.
The correction is Gaussian divergence-free, so \cref{app:eq:hermite-div-free} gives
\begin{equation*}
    \mbz^\top\mbnu^{(\ell)}(\mbz,t)=\nabla_{\mbz}\cdot\mbnu^{(\ell)}(\mbz,t).
\end{equation*}

For $\ell=0$, the correction is zero.  For $1\leq\ell\leq3$, the right-hand side is a polynomial of degree at most $\ell-1$, and hence for some constant $C > 0$
\begin{equation*}
    \left|\mbz^\top\mbnu^{(\ell)}(\mbz,t)\right| \leq CV(\mbz), \qquad (\mbz,t) \in \bbR^K \times [0,T],
\end{equation*}
where $V$ is the Lyapunov function $V(\mbz)=1+\|\mbz\|^2$. 
The infinitesimal generator of \cref{app:eq:whitened-posterior-sde} satisfies
\begin{equation*}
    \cL_t V(\mbz)=-\mbz^\top\mbC(t)\mbz +2\mbz^\top\mbnu^{(\ell)}(\mbz,t) +\mathrm{tr}(\mbC(t)),
\end{equation*}
and is therefore bounded above by $C V(\mbz)$.  Moreover,
\begin{equation*}
    \lim_{R\to\infty} \inf_{\|\mbz\|>R}V(\mbz)=\infty.
\end{equation*}

The non-explosion criterion from \citet[Theorem~3.5]{khasminskii2011stochastic} therefore implies that a pathwise-unique, strong solution is defined on all of $[0, T]$. 
This proves \Cref{app:assum:regularity} for the posterior SDE.
Notice that our argument does not require the posterior drift to have linear
growth.

\paragraph{Fokker--Planck equation and Gaussian marginals.}
By construction, $\{q(\mbx,t)\}_{0 \leq t \leq T}$ is a weak solution to the Fokker--Planck equation for the reference drift $\mbf_q$.  
Moreover,
\begin{equation*}
    \nabla_{\mbx}\cdot  \left[q(\mbx,t)\mbh^{(\ell)}(\mbx,t)\right] =0, \quad (\mbx, t) \in \bbR^K \times [0, T]
\end{equation*}
so the same $q(\mbx,t)$ also solves the weak Fokker--Planck equation for the corrected drift $\mbf_q+\mbh^{(\ell)}$.  
The corrected drift is polynomial, $\mbSigma(t)$ is bounded, and Gaussian distributions have moments of every finite order.  
Hence $q$ belongs to the regularity class \cref{app:eqn:fpk-regularity-class}.

It remains to show that, within the regularity class \cref{app:eqn:fpk-regularity-class}, the Fokker--Planck equation has a unique solution. 
The superposition principle (\citet[Theorem~2.5]{trevisan2016well}) states that any weak Fokker--Planck solution $\{\rho(\mbx, t)\}_{0 \leq t \leq T}$ can be lifted to a probability distribution over continuous paths whose one-time marginals are $\{\rho(\mbx, t)\}_{0 \leq t \leq T}$ and which is a weak solution to the SDE with drift $\mbf_q+\mbh^{(\ell)}$ and diffusion coefficient $\mbSigma^{1/2}$.

Suppose $\{q(\mbx, t)\}_{0 \leq t \leq T}$ and $\{\tilde{q}(\mbx, t)\}_{0 \leq t \leq T}$ are two weak solutions to the Fokker--Planck equation with drift $\mbf_q + \mbh^{(\ell)}$ and diffusion coefficient $\mbSigma^{1/2}$. 
By the superposition principle, we obtain two probability distributions having one-time marginals $\{q(\mbx, t)\}_{0 \leq t \leq T}$ and $\{\tilde{q}(\mbx, t)\}_{0 \leq t \leq T}$ that are weak solutions to the variational SDE  with drift $\mbf_q + \mbh^{(\ell)}$ and diffusion coefficient $\mbSigma^{1/2}$.
The strong existence and pathwise uniqueness established above imply that this SDE is unique in law.
Therefore, the two weak solutions have the same path law and hence the same one-time marginals. 
Thus $q(\mbx,t)=\tilde q(\mbx,t)$ almost everywhere in $\mbx$ for every $t\in[0,T]$.
This verifies \Cref{app:ass:fpk-unique} for the posterior drift.

\paragraph{Densities.}
The posterior marginals are non-degenerate Gaussians and hence have smooth, strictly positive Lebesgue densities. 
This verifies the posterior part of
\Cref{app:assum:densities}.

\paragraph{Absolute continuity.}
Define the first exit time from the ball of radius $n$
\begin{equation*}
    \tau_n  = \inf\left\{t\in[0,T] \ | \  \|\mbx(t)\|_2\geq n \right\}, \qquad \inf\varnothing=\infty.
\end{equation*}
Since the prior and posterior drifts are locally bounded and $\mbSigma(t)^{-1/2}$ is bounded on $[0,T]$, there exists $K_n<\infty$ such that
\begin{equation*}
    \sup_{\substack{\|\mby\|_2\leq n\\0\leq t\leq T}} \left\| \mbSigma(t)^{-1/2} \left[ \mbf_q(\mby,t)+\mbh^{(\ell)}(\mby,t) -\mbf_p(\mby,t) \right] \right\|_2 \leq K_n.
\end{equation*}
Consequently,
{\small
\begin{equation*}
\begin{aligned}
    \bbE_{\bbP}\left[ \exp\left\{ \frac{1}{2} \int_0^{T\wedge\tau_n} \left\| \mbSigma(t)^{-1/2} \left[ \mbf_q(\mbx(t),t) +\mbh^{(\ell)}(\mbx(t),t) -\mbf_p(\mbx(t),t) \right] \right\|_2^2dt \right\} \right]  \leq  \exp\left\{\frac{1}{2}K_n^2T\right\} <\infty.
\end{aligned}
\end{equation*}
}
Thus, Novikov's criterion (\citet[Corollary~3.5.13]{karatzas2014brownian}) holds for the stopped process.

Both initial densities are non-degenerate Gaussian densities and are therefore strictly positive on $\bbR^K$. Define
{\small
\begin{equation*}
\begin{aligned}
    \cZ_T^{(n)}
    &= \frac{q(\mbx(0),0)}{p(\mbx(0),0)} \exp\Bigg\{ \sum_{i=1}^K \int_0^{T\wedge\tau_n} \left[ \mbSigma(t)^{-1/2} \left\{ \mbf_q(\mbx(t),t) +\mbh^{(\ell)}(\mbx(t),t) -\mbf_p(\mbx(t),t) \right\} \right]_i d\mbw_i(t)\\
    &-\frac{1}{2} \int_0^{T\wedge\tau_n} \left\| \mbSigma(t)^{-1/2} \left[ \mbf_q(\mbx(t),t)  +\mbh^{(\ell)}(\mbx(t),t)-\mbf_p(\mbx(t),t) \right] \right\|_2^2dt\Bigg\},
\end{aligned}
\end{equation*}
}
where $(\mbw(t))_{0\leq t\leq T}$ is a $K$-dimensional Brownian motion with respect to $\bbP$.

Novikov's criterion implies that the stochastic exponential appearing in $\cZ_T^{(n)}$ is a martingale. Conditionally on $\mbx(0)$, it has expectation one, and therefore
\begin{equation*}
    \bbE_{\bbP}\left[\cZ_T^{(n)}\right]
    = \bbE_{\bbP}\left[ \frac{q(\mbx(0),0)}{p(\mbx(0),0)}   \right] = \int \frac{q(\mbx,0)}{p(\mbx,0)} p(\mbx,0)\,d\mbx =1.
\end{equation*}
We may therefore define a probability measure
$\widehat{\bbQ}^{(n)}$ on $\mathcal F_T$ by
\begin{equation*}
    \frac{d\widehat{\bbQ}^{(n)}}{d\bbP} = \cZ_T^{(n)}.
\end{equation*}
Under $\widehat{\bbQ}^{(n)}$, Girsanov's theorem shows that
\begin{equation*}
\begin{aligned}
    \widehat{\mbw}^{(n)}(t) &= \mbw(t) - \int_0^{t\wedge\tau_n} \mbSigma(s)^{-1/2} \left\{ \mbf_q(\mbx(s),s) +\mbh^{(\ell)}(\mbx(s),s) -\mbf_p(\mbx(s),s) \right\}ds
\end{aligned}
\end{equation*}
is a Brownian motion.
Hence, up to $\tau_n$, the coordinate process satisfies
\begin{equation*}
    d\mbx(t) = \left[ \mbf_q(\mbx(t),t) +\mbh^{(\ell)}(\mbx(t),t) \right]dt + \mbSigma(t)^{1/2} d\widehat{\mbw}^{(n)}(t),
\end{equation*}
and its initial density is $q(\mbx,0)$.

Let $\cG_n = \sigma\left\{ \mbx(t\wedge\tau_n):0\leq t\leq T \right\}$ be the sigma-algebra generated by the stopped coordinate process.
Under both $\widehat{\bbQ}^{(n)}$ and $\bbQ$, the stopped coordinate process has initial density $q(\mbx,0)$ and satisfies the same stopped posterior SDE. 
Uniqueness in law therefore gives
\begin{equation*}
    \left.
    \widehat{\bbQ}^{(n)}
    \right|_{\mathcal G_n}
    =
    \left.
    \bbQ
    \right|_{\mathcal G_n}.
\end{equation*}
Because $\mbSigma(t)$ is invertible, the prior Brownian motion up to
$\tau_n$ is determined by the stopped coordinate process. Hence
$\cZ_T^{(n)}$ is $\mathcal G_n$-measurable. It follows that, for every
$A\in\mathcal G_n$,
\begin{equation*}
    \bbQ(A)
    =
    \widehat{\bbQ}^{(n)}(A)
    =
    \bbE_{\bbP}\left[
        \mathbbm{1}_A\cZ_T^{(n)}
    \right].
\end{equation*}

Define the corresponding unstopped stochastic exponential by
{\small
\begin{equation*}
\begin{aligned}
    \cZ_T &= \frac{q(\mbx(0),0)}{p(\mbx(0),0)} \exp\Bigg\{ \sum_{i=1}^K \int_0^T \left[ \mbSigma(t)^{-1/2} \left\{ \mbf_q(\mbx(t),t) +\mbh^{(\ell)}(\mbx(t),t) -\mbf_p(\mbx(t),t) \right\} \right]_i d\mbw_i(t) \\
    &- \frac{1}{2} \int_0^T \left\|    \mbSigma(t)^{-1/2} \left[\mbf_q(\mbx(t),t) +\mbh^{(\ell)}(\mbx(t),t)-\mbf_p(\mbx(t),t) \right] \right\|_2^2dt\Bigg\}.
\end{aligned}
\end{equation*}
}
This random variable is well defined $\bbP$-almost surely by local
boundedness and non-explosion.

To remove the stopping time, let $E_n=\{\tau_n>T\}$. 
Since the prior and posterior processes have strong solutions defined on all of $[0,T]$, their paths are continuous and bounded on $[0,T]$. 
Therefore, $E_n\uparrow\Omega$ both $\bbP$-almost surely and $\bbQ$-almost surely. 
Moreover, by definition, $\cZ_T^{(n)}=\cZ_T$ on $E_n$. 
For every $A$ measurable with respect to $\sigma(\mbx(t) \ | \ 0 \leq t \leq T)$, the event $A\cap E_n$ belongs to $\cG_n$, and hence
\begin{equation*}
\begin{aligned}
    \bbQ(A\cap E_n) &= \bbE_{\bbP}\left[ \mathbbm{1}_{A\cap E_n}\cZ_T^{(n)}  \right] = \bbE_{\bbP}\left[\mathbbm{1}_{A\cap E_n}\cZ_T \right].
\end{aligned}
\end{equation*}
Continuity from below under $\bbQ$ and monotone convergence under
$\bbP$ now give
\begin{equation*}
    \bbQ(A) = \lim_{n\to\infty}\bbQ(A\cap E_n) = \lim_{n\to\infty} \bbE_{\bbP}\left[\mathbbm{1}_{A\cap E_n}\cZ_T \right] = \bbE_{\bbP}\left[ \mathbbm{1}_A\cZ_T \right].
\end{equation*}
Therefore,
\begin{equation*}
    \frac{d\bbQ}{d\bbP}  = \cZ_T.
\end{equation*}
In particular, $\bbQ$ is absolutely continuous with respect to $\bbP$. 
Since $\cZ_T>0$ $\bbP$-almost surely, the reverse absolute continuity also holds. 
Therefore, the prior and posterior path laws are mutually absolutely continuous.

\paragraph{Finite KL divergence.}
The initial-distribution term is finite because both initial distributions are non-degenerate Gaussians.

For the path-space term, observe that
\begin{equation*}
    \mbf_p(\mbx,t) -\left\{\mbf_q(\mbx,t)+\mbh^{(\ell)}(\mbx,t)\right\}  =\mbr(\mbx,t)-\mbh^{(\ell)}(\mbx,t).
\end{equation*}

Both terms on the right have polynomial growth in $\mbx$, uniformly for $t\in[0,T]$, while $\mbSigma(t)^{-1}$ is uniformly bounded.
Since the posterior marginals are Gaussian and their means and covariances are bounded on $[0,T]$, they have uniformly bounded moments of every finite order.
Consequently,
\begin{equation*}
    \int_0^T \bbE_{\bbQ} \left\|    \mbSigma(t)^{-1/2} \left[ \mbf_p(\mbx(t),t) -\left\{\mbf_q(\mbx(t),t) +\mbh^{(\ell)}(\mbx(t),t)\right\}  \right] \right\|^2dt <\infty.
\end{equation*}
Thus \Cref{app:ass:finite-kl} holds, and the path-space KL identity \cref{app:eqn:path-kl} is finite.

We have therefore verified the assumptions of \Cref{app:subsec:assumptions} for the Helmholtz-SDE variational family with $0 \leq \ell \leq 3$, conditional on the stated assumptions on the prior.  
For $\ell \geq 4$, the divergence-free identity still preserves the prescribed Fokker--Planck equation, but the Lyapunov estimate above does not establish non-explosion. 
Global well-posedness must be verified separately.

\subsection{Optimality of Helmholtz-SDE}
\label{app:subsec:optimality}

Now that we have established that Helmholtz-SDE produces a valid variational SDE, we study its optimality.
First, we prove \Cref{thm:helm-optimal}, which shows Helmholtz-SDE approximates the posterior arbitrarily closely in the conjugate setting.
We then study the nonconjugate setting. 
We quantify the suboptimality resulting from restricting the $q$-divergence-free correction to a polynomial vector field of degree $\ell$.
As a result, we get a lower bound on the improvement over the variational family defined by the reference drift.

\subsubsection{Conjugate setting}
\label{app:subsec:optimality-conj}

\begin{proof}[Proof of \Cref{thm:helm-optimal}]
Suppose that the prior SDE is linear and that the observation model is linear and Gaussian.
The posterior path law $\bbQ^\star$ is then Gauss--Markov, its drift $\mbf^\star(\mbx,t)$ is affine in $\mbx$, and its one-time marginals have the form $q^\star(\mbx,t)=\cN(\mbm^\star(t),\mbS^\star(t))$ \citep{kalman1961new,sarkka2023bayesian}.
The covariance $\mbS^\star(t)$ is positive definite.
The functions $\mbm^\star(t)$ and $\mbS^\star(t)$ are continuous and piecewise continuously differentiable, with possible derivative jumps at the observation times \citep[Appendix J]{hu2025sing}.

For any continuously differentiable mean and covariance paths $(\mbm(t),\mbS(t))_{0\leq t\leq T}$, Helmholtz-SDE uses the affine reference drift from \Cref{app:prop:gaussian-compatible-drift} and adds the degree-$1$ correction obtained from \Cref{app:prop:hermite-projection}.
Because the prior and reference drifts are affine, the residual vector field is affine.
Moreover, the identity \cref{app:eq:hermite-raising} shows that the divergence-free projection preserves the Hermite degree.
Consequently, the projection onto the subspace $\cV_{t, \ell}$ of $q$-divergence-free vector fields that are degree $\ell$ polynomials is equal to the projection onto the subspace $\cV_t$ of all $q$-divergence-free vector fields.
Thus, at any prescribed Gaussian marginals, Helmholtz-SDE selects the drift that minimizes the path-space KL divergence to the prior among all admissible drifts having those marginals.

The paths $\mbm^\star$ and $\mbS^\star$ need not be continuously differentiable at the observation times, so $\bbQ^\star$ need not itself belong to the variational family $\cQ_{(1)}$.
We instead argue that $\bbQ^\star$ can be approximated arbitrarily closely in KL by $\cQ_{(1)}$.
Fix $\delta>0$, and let $U_\delta\subset[0,T]$ be a union of neighborhoods of the observation times with Lebesgue measure at most $\delta$.
The mean and covariance paths can be smoothed inside $U_\delta$ to obtain continuously differentiable functions $\mbm^\delta$ and $\mbS^\delta$ that agree with $\mbm^\star$ and $\mbS^\star$ outside $U_\delta$ and converge uniformly to them as $\delta \to 0$.
The smoothing may be chosen so that the first derivatives remain uniformly bounded in $\delta$.
For sufficiently small $\delta$, $\mbS^\delta(t)$ remains positive definite because $\lambda_{\min}(\mbS^\star(t))$ is bounded away from zero on $[0,T]$.

Let $\bbQ^\delta$ be the path law produced by degree-$1$ Helmholtz-SDE applied to $\mbm^\delta$ and $\mbS^\delta$.
Outside $U_\delta$, these marginal paths and their derivatives agree with those of $\bbQ^\star$.
The corresponding Helmholtz-SDE drift therefore minimizes the fixed-marginal path-space KL objective at $q^\star$.
The posterior drift $\mbf^\star$ minimizes the same objective, since replacing it by a drift with the same marginals would leave the expected log likelihood and the initial KL term unchanged while changing only the path-space KL term.
Strict convexity of the fixed-marginal objective gives
\begin{equation*}
\mbf_q^\delta(\mbx,t)+\mbh^{(1),\delta}(\mbx,t)=\mbf^\star(\mbx,t), \qquad t\notin U_\delta,
\end{equation*}
up to $q^\star(\mbx,t)$-almost-everywhere equality.

The drifts $\mbf_q^\delta + \mbh^{(1),\delta}$ and $\mbf^\star$ are affine in $\mbx$.
The uniform bounds on the smoothed marginal paths and their derivatives imply that their affine coefficients are uniformly bounded on $U_\delta$.
The Gaussian marginals of $\bbQ^\delta$ also have uniformly bounded second moments.
Since $\mbSigma(t)^{-1}$ is bounded on $[0,T]$, there is a constant $C<\infty$, independent of $\delta$, such that
\begin{equation*}
\bbE_{\bbQ^\delta}\left\|\mbSigma(t)^{-1/2}\left\{\mbf_q^\delta(\mbx(t),t)+\mbh^{(1),\delta}(\mbx(t),t)-\mbf^\star(\mbx(t),t)\right\}\right\|_2^2\leq C, \qquad t\in U_\delta.
\end{equation*}

Applying the path-space KL identity to $\bbQ^\delta$ and $\bbQ^\star$ gives
{\small
\begin{equation*}
\begin{aligned}
&\KL{\bbQ^\delta}{\bbQ^\star}\\
&=\KL{q^\delta(\mbx,0)}{q^\star(\mbx,0)}+\frac{1}{2}\int_{U_\delta}\bbE_{\bbQ^\delta}\left\|\mbSigma(t)^{-1/2}\left\{\mbf_q^\delta(\mbx(t),t)+\mbh^{(1),\delta}(\mbx(t),t)-\mbf^\star(\mbx(t),t)\right\}\right\|_2^2dt\\
&\leq\KL{q^\delta(\mbx,0)}{q^\star(\mbx,0)}+\frac{C\delta}{2}.
\end{aligned}
\end{equation*}
}
The Gaussian initial KL converges to zero by uniform convergence of the mean and covariance paths.
Hence $\KL{\bbQ^\delta}{\bbQ^\star}\to0$ as $\delta\to0$, and therefore $\inf_{\bbQ\in\cQ_{(1)}}\KL{\bbQ}{\bbQ^\star}=0$.
\end{proof}

\subsubsection{Nonconjugate setting}
\label{app:subsec:optimality-nonconj}

In the nonconjugate setting, the posterior generally does not have Gaussian one-time marginals.
Simulation-free variational inference algorithms that impose Gaussian one-time marginals therefore share a common, irreducible approximation error. 
Helmholtz-SDE further restricts the variational family by considering a subset of divergence-free vector fields.

As in \Cref{app:thm:hilbert-projection}, we consider the Hilbert space $\cL^2(q(\mbx, t), \bbR^K, \mbSigma(t))^K$ of vector fields satisfying $\bbE_{q(\mbx, t)}\|\mbf(\mbx)\|^2 < \infty$ equipped with the inner product defined in \cref{app:eqn:l2-inner-prod}.
Let $\cV_t$ be the closed subspace of weakly $q$-divergence-free vector fields, let $\cP_\ell$ be the subspace of polynomial vector fields of degree at most $\ell$, and let $\cV_{t,\ell}=\cV_t\cap\cP_\ell$.

Let $\Pi_\cS$ be the orthogonal projection onto the subspace $\cS$.
Then the degree $\ell$ correction is $\mbh^{(\ell)}(\mbx,t)= [\Pi_{\cV_{t,\ell}}\mbr(\cdot,t)](\mbx)$.
First, we quantify the reduction in negative ELBO at \emph{fixed marginals} from increasing the degree of the polynomial correction.
\begin{proposition}[Improvement at fixed marginals]
\label{app:prop:polynomial-suboptimality}
Suppose \Cref{app:assum:resid} holds, and fix one-time marginals $\{q(\mbx, t)\}_{0 \leq t \leq T}$.
Let $\ell \geq 0$ and $k \geq 1$ be integers for which the degree $\ell$ and $\ell+k$ corrections induce admissible path laws.

The reduction in negative ELBO from using the degree $\ell+k$ correction in place of the degree $\ell$ correction is
\begin{equation}\label{app:eq:poly-suboptimality}
\begin{aligned}
\Delta_{\ell,\ell+k}(q)&=\frac{1}{2}\int_0^T \left\{\|\mbr(\mbx,t)-\mbh^{(\ell)}(\mbx,t)\|_{q,t}^2-\|\mbr(\mbx,t)-\mbh^{(\ell+k)}(\mbx,t)\|_{q,t}^2\right\}dt\\
&=\frac{1}{2}\int_0^T\|\mbh^{(\ell+k)}(\mbx,t)-\mbh^{(\ell)}(\mbx,t)\|_{q,t}^2dt.
\end{aligned}
\end{equation}
Equivalently,
\begin{equation}\label{app:eq:degree-est}
\Delta_{\ell,\ell+k}(q)=\frac{1}{2}\sum_{\ell<|\mbalpha|\leq\ell+k}\int_0^T\widetilde{\mbc}_{\mbalpha}(t)^{\top}\mbLambda(t)^{-1}\widetilde{\mbc}_{\mbalpha}(t)dt.
\end{equation}
$\widetilde{\mbc}_{\mbalpha}(t)$ are the coefficients of the $q$-divergence-free vector field in the Hermite basis obtained from \Cref{app:prop:hermite-projection}.
\end{proposition}

\begin{proof}
Since the path laws obtained from $\mbh^{(\ell)}$ and $\mbh^{(\ell+k)}$ possess one-time marginals $\{q(\mbx, t)\}_{0 \leq t \leq T}$, their reconstruction losses and initial KL divergences are equal.
Therefore, the reduction in negative ELBO is the difference between the path-space KL divergences to the prior.
The second equality in \cref{app:eq:poly-suboptimality} follows from the Pythagorean identity for projections onto the nested subspaces $\cV_{t,\ell}\subseteq\cV_{t,\ell+k}$.

The identity \cref{app:eq:hermite-raising} maps each Hermite degree to the next degree and does not couple different degrees.
Consequently, $\cV_t$ is the orthogonal direct sum of its Hermite-degree components, and the degree $\ell$ and $\ell+k$ corrections have identical Hermite coefficients through degree $\ell$.
Orthonormality of the Hermite basis gives \cref{app:eq:degree-est}.
\end{proof}

\Cref{app:prop:polynomial-suboptimality} compares polynomial degrees at fixed marginals. 
We now consider optimizing within each variational family.
Define $\cQ_{(\ell)}$ to be the variational family obtained by applying Helmholtz-SDE to the Gaussian one-time marginals whose mean and covariance satisfy \Cref{app:assum:resid}.
From \cref{app:eq:degree-est}, we obtain a lower bound on the improvement from using $\cQ_{(\ell + k)}$ over $\cQ_{(\ell)}$.
\begin{proposition}[Suboptimality of variational families]
\label{app:prop:error-nonlinear}
Let $\ell \geq 0$ and $k \geq 1$ be integers for which the degree $\ell$ and $\ell+k$ corrections induce admissible path laws.
For any $\epsilon>0$, suppose that $\bbQ_\epsilon\in\cQ_{(\ell)}$ has marginals $q_\epsilon(\mbx,t)$ and satisfies
\begin{equation*}
\KL{\bbQ_\epsilon}{\bbQ^\star}\leq\inf_{\bbQ\in\cQ_{(\ell)}}\KL{\bbQ}{\bbQ^\star}+\epsilon.
\end{equation*}
Let $\mbh_\epsilon^{(\ell)}$ and $\mbh_\epsilon^{(\ell+k)}$ be the degree $\ell$ and $\ell + k$  corrections computed at $q_\epsilon(\mbx,t)$.
Then
\begin{equation}\label{app:eq:optimality-bound}
\begin{aligned}
&\inf_{\bbQ\in\cQ_{(\ell)}}\KL{\bbQ}{\bbQ^\star}-\inf_{\bbQ\in\cQ_{(\ell + k)}}\KL{\bbQ}{\bbQ^\star}\\
&\geq\frac{1}{2}\int_0^T\|\mbh_\epsilon^{(\ell + k)}(\mbx,t) - \mbh_\epsilon^{(\ell)}(\mbx,t)\|_{q_\epsilon,t}^2dt-\epsilon.
\end{aligned}
\end{equation}
\end{proposition}

\begin{proof}
Let $\bbQ_\epsilon^{(\ell + k)}$ be the path law obtained by adding $\mbh_\epsilon^{(\ell + k)}$ to the reference drift of $\bbQ_\epsilon$.
Then $\bbQ_\epsilon^{(\ell + k)}$ and $\bbQ_\epsilon$ have the same one-time marginals because $\mbh_\epsilon^{(\ell + k)}$ is $q_\epsilon$-divergence-free.
Their reconstruction losses and initial KL divergences are equal.
The difference between their KL divergences to $\bbQ^\star$ is consequently the difference between their path-space KL divergences to the prior.
By \Cref{app:prop:polynomial-suboptimality} and the $\epsilon$-optimality of $\bbQ_\epsilon$,
\begin{equation*}
\begin{aligned}
&\inf_{\bbQ\in\cQ_{(\ell)}}\KL{\bbQ}{\bbQ^\star}-\inf_{\bbQ\in\cQ_{(\ell + k)}}\KL{\bbQ}{\bbQ^\star}\\
&\geq\KL{\bbQ_\epsilon}{\bbQ^\star}-\KL{\bbQ_\epsilon^{(\ell + k)}}{\bbQ^\star}-\epsilon\\
&= \Delta_{\ell,\ell+k}(q_\epsilon) -\epsilon\\
&=\frac{1}{2}\int_0^T\|\mbh_\epsilon^{(\ell + k)}(\mbx,t) - \mbh_\epsilon^{(\ell)}(\mbx,t)\|_{q_\epsilon,t}^2dt-\epsilon
\end{aligned}
\end{equation*}
\end{proof}
In particular, by choosing $\ell = 0$, \Cref{app:prop:error-nonlinear} provides a lower bound on the improvement over the variational family defined by the reference drift (e.g., SVISE \cite{course2023state,course2023amortized} and SDE Matching \cite{bartosh2025sde}).


\Cref{app:prop:error-nonlinear} provides a heuristic for estimating the improvement of $\cQ_{(\ell + k)}$ over $\cQ_{(\ell)}$.
After fitting Helmholtz-SDE with a degree $\ell$ correction, we hold the fitted marginals and generative model parameters fixed, compute the Hermite coefficients through degree $\ell+k$, and evaluate \cref{app:eq:optimality-bound}.
We use the resulting
quantity as an estimate of the improvement, taking the unknown
optimization error $\epsilon$ to be zero.
With an optimal degree $\ell$ fit and exact coefficient evaluation, this quantity is a lower bound on the improvement between the  variational families.
We validate this heuristic for the stochastic Lorenz attractor in \Cref{app:subsec:lorenz}.

\subsection{Algorithm summary}
\label{app:subsec:alg-summary}

We summarize the Hermite implementation of the polynomial correction in \Cref{app:alg:helm_correction} and the resulting Helmholtz-SDE training procedure in \Cref{app:alg:helmholtz_sde}.
The coefficient samples used by \texttt{HelmCorrection} are independent of the state samples used to estimate the path-space KL.

\clearpage
\begin{algorithm}[H]
\caption{\texttt{HelmCorrection}: Hermite approximation to the divergence-free correction}
\label{app:alg:helm_correction}
\Require{polynomial order $\ell$; prior drift $\mbf_p$; reference drift $\mbf_q$; diffusion covariance $\mbSigma$; posterior mean $\mbm$ and covariance $\mbS$ at time $t$; number of coefficient samples $N_{\mathrm{mc}}$}

\medskip
\tcp{Whiten and rotate}
\AlgStep{Compute the symmetric square root $\mbS^{1/2}$ and its inverse $\mbS^{-1/2}$}

\smallskip
\AlgStep{Form $\mbC\leftarrow\mbS^{-1/2}\mbSigma(t)\mbS^{-1/2}$ and compute $\mbC=\mbQ\mbLambda\mbQ^\top$ with $\mbLambda=\diag(\lambda_1,\ldots,\lambda_K)$}

\medskip
\tcp{Estimate the Hermite coefficients of the rotated residual}
\AlgStep{Sample $\mbv^{(n)}\overset{\mathrm{i.i.d.}}{\sim}\cN(\mb0,\bbI_K)$ for $n=1,\ldots,N_{\mathrm{mc}}$}

\smallskip
\AlgStep{Set $\mbx^{(n)}\leftarrow\mbm+\mbS^{1/2}\mbQ\mbv^{(n)}$}

\smallskip
\AlgStep{Compute $\widetilde{\mbrho}^{(n)}\leftarrow\mbQ^\top\mbS^{-1/2}\left\{\mbf_p(\mbx^{(n)},t)-\mbf_q(\mbx^{(n)},t)\right\}$}

\smallskip
\For{$\mbalpha\in\cN_{K,\ell}$ and $i=1,\ldots,K$}{
    \AlgStep{$\widehat{\widetilde{\mbb}}_{\mbalpha,i}\leftarrow N_{\mathrm{mc}}^{-1}\sum_{n=1}^{N_{\mathrm{mc}}}\widetilde{\mbrho}^{(n)}_i\hat{H}_{\mbalpha}(\mbv^{(n)})$}
}

\medskip
\tcp{Project coefficient by coefficient onto the divergence-free subspace}
\For{$\mbbeta\in\cN_{K,\ell+1}\setminus\cN_{K,0}$}{
    \AlgStep{$\widehat\mu_{\mbbeta}\leftarrow\left\{\sum_{i=1}^K\sqrt{\mbbeta_i}\widehat{\widetilde{\mbb}}_{\mbbeta-\mbe_i,i}\right\}\big/\left\{\sum_{i=1}^K\lambda_i\mbbeta_i\right\}$}
}

\smallskip
\For{$\mbalpha\in\cN_{K,\ell}$ and $i=1,\ldots,K$}{
    \AlgStep{$\widehat{\widetilde{\mbc}}_{\mbalpha,i}\leftarrow\widehat{\widetilde{\mbb}}_{\mbalpha,i}-\lambda_i\sqrt{\mbalpha_i+1}\widehat\mu_{\mbalpha+\mbe_i}$}
}

\medskip
\tcp{Return the correction in the original coordinates}
\AlgStep{Define $\widehat{\mbnu}^{(\ell)}(\mbv,t)\leftarrow\sum_{\mbalpha\in\cN_{K,\ell}}\widehat{\widetilde{\mbc}}_{\mbalpha}\hat{H}_{\mbalpha}(\mbv)$}

\smallskip
\AlgStep{Define $\widehat{\mbh}^{(\ell)}(\mbx,t)\leftarrow\mbS^{1/2}\mbQ\widehat{\mbnu}^{(\ell)}\left(\mbQ^\top\mbS^{-1/2}(\mbx-\mbm),t\right)$}

\smallskip
\Return $\widehat{\mbh}^{(\ell)}(\mbx,t)$
\end{algorithm}


\begin{algorithm}[H]
\caption{Helmholtz-SDE training}
\label{app:alg:helmholtz_sde}
\Require{dataset $\cD=\{\mby^{(m)}(t_n)\}_{n=1,m=1}^{N,M}$; prior drift $\mbf_p(\mbx,t\mid\mbtheta)$; diffusion covariance $\mbSigma(t\mid\mbtheta)$; observation model $p(\mby(t_n)\mid\mbx(t_n),\mbtheta)$; model parameters $\mbtheta$; parameters $\mbphi$ of the posterior marginal paths $(\mbm^{(m)}(t\mid\mbphi),\mbS^{(m)}(t\mid\mbphi))_{0\leq t\leq T}$; reference drift $\mbf_q$; polynomial order $\ell$; sample sizes $N_{\mathrm{time}}$, $N_{\mathrm{space}}$, and $N_{\mathrm{mc}}$}

\For{learning iterations}{
    \tcp{Compute the initial KL term}
    \AlgStep{$\cL_{\mathrm{init}}\leftarrow\sum_{m=1}^M\KL{\cN(\mbm^{(m)}(0\mid\mbphi),\mbS^{(m)}(0\mid\mbphi))}{\cN(\mbmu_0,\mbV_0)}$}

    \medskip
    \tcp{Estimate the reconstruction term}
    \AlgStep{Sample $s_j\overset{\mathrm{i.i.d.}}{\sim}\mathrm{Unif}\{t_1,\ldots,t_N\}$ for $j=1,\ldots,N_{\mathrm{time}}$}

    \smallskip
    \AlgStep{Sample $\mbx_{m,j,k}\overset{\mathrm{i.i.d.}}{\sim}\cN(\mbm^{(m)}(s_j\mid\mbphi),\mbS^{(m)}(s_j\mid\mbphi))$ for $k=1,\ldots,N_{\mathrm{space}}$}

    \smallskip
    \AlgStep{$\widehat{\cL}_{\mathrm{rec}}\leftarrow-\frac{N}{N_{\mathrm{time}}N_{\mathrm{space}}}\sum_{m=1}^M\sum_{j=1}^{N_{\mathrm{time}}}\sum_{k=1}^{N_{\mathrm{space}}}\log p(\mby^{(m)}(s_j)\mid\mbx_{m,j,k},\mbtheta)$}

    \medskip
    \tcp{Estimate the path-space KL term}
    \AlgStep{Sample $s_j\overset{\mathrm{i.i.d.}}{\sim}\mathrm{Unif}[0,T]$ for $j=1,\ldots,N_{\mathrm{time}}$}

    \smallskip
    \For{$m=1,\ldots,M$ and $j=1,\ldots,N_{\mathrm{time}}$}{
        \AlgStep{$\widehat{\mbh}_{m,j}^{(\ell)}(\cdot)\leftarrow\text{\texttt{HelmCorrection}}(\ell,\mbf_p,\mbf_q,\mbSigma,\mbm^{(m)}(s_j\mid\mbphi),\mbS^{(m)}(s_j\mid\mbphi),s_j,N_{\mathrm{mc}})$}
    }

    \smallskip
    \AlgStep{Independently sample $\mbx_{m,j,k}\overset{\mathrm{i.i.d.}}{\sim}\cN(\mbm^{(m)}(s_j\mid\mbphi),\mbS^{(m)}(s_j\mid\mbphi))$ for $k=1,\ldots,N_{\mathrm{space}}$}

    \smallskip
    \AlgStep{$\widehat{\cL}_{\mathrm{path}}\leftarrow\frac{T}{2N_{\mathrm{time}}N_{\mathrm{space}}}\sum_{m=1}^M\sum_{j=1}^{N_{\mathrm{time}}}\sum_{k=1}^{N_{\mathrm{space}}}\left\|\mbSigma(s_j\mid\mbtheta)^{-1/2}\left\{\mbf_p(\mbx_{m,j,k},s_j\mid\mbtheta)-\mbf_q(\mbx_{m,j,k},s_j)-\widehat{\mbh}_{m,j}^{(\ell)}(\mbx_{m,j,k})\right\}\right\|_2^2$}

    \medskip
    \tcp{Update the model and variational parameters}
    \AlgStep{$\cL\leftarrow\widehat{\cL}_{\mathrm{rec}}+\cL_{\mathrm{init}}+\widehat{\cL}_{\mathrm{path}}$}

    \smallskip
    \AlgStep{$\mbtheta,\mbphi\leftarrow\texttt{gradient-step}(\cL,\mbtheta,\mbphi)$}
}
\end{algorithm}
\clearpage